% =============================================================================
% LUCRARE DE LICENȚĂ — PrimeRent
% Universitatea Transilvania din Brașov
% Facultatea de Matematică și Informatică — Domeniul INFORMATICĂ
% =============================================================================
\documentclass[12pt,a4paper,oneside]{report}

% ─── Fonturi și encoding ─────────────────────────────────────────────────────
\usepackage[T1]{fontenc}
\usepackage[utf8]{inputenc}
\usepackage[romanian]{babel}
% Font Calibri conform ghidului — se folosește helvet (Helvetica/Arial) ca alternativă
% compatibilă pe toate distribuțiile LaTeX. Pentru Calibri nativ,
% compilați cu XeLaTeX/LuaLaTeX și \usepackage{fontspec}\setmainfont{Calibri}
\usepackage{helvet}
\renewcommand{\familydefault}{\sfdefault}

% ─── Geometrie pagină ────────────────────────────────────────────────────────
\usepackage[
  left=3cm,
  right=2cm,
  top=2.5cm,
  bottom=2.5cm
]{geometry}

% ─── Spațiere la un rând ─────────────────────────────────────────────────────
\usepackage{setspace}
\singlespacing

% ─── Pachete generale ────────────────────────────────────────────────────────
\usepackage{graphicx}
\usepackage{float}
\usepackage{booktabs}
\usepackage{longtable}
\usepackage{array}
\usepackage{multirow}
\usepackage{enumitem}
\usepackage{xcolor}
\usepackage{hyperref}
\usepackage{url}
\usepackage{caption}
\usepackage{subcaption}
\usepackage{fancyhdr}
\usepackage{titlesec}
\usepackage{tocloft}
\usepackage{amsmath}
\usepackage{amssymb}
\usepackage{pifont}

% ─── Cod sursă ───────────────────────────────────────────────────────────────
\usepackage{listings}
\usepackage{inconsolata}

\lstset{
  basicstyle=\footnotesize\ttfamily,
  breaklines=true,
  breakatwhitespace=false,
  columns=fullflexible,
  keepspaces=true,
  frame=single,
  framesep=5pt,
  framerule=0.5pt,
  rulecolor=\color{gray!50},
  backgroundcolor=\color{gray!5},
  numbers=left,
  numberstyle=\tiny\color{gray},
  numbersep=8pt,
  showstringspaces=false,
  tabsize=2,
  captionpos=b,
  xleftmargin=15pt,
  xrightmargin=5pt,
  aboveskip=10pt,
  belowskip=10pt,
  literate={ă}{{\u{a}}}1 {î}{{\^{i}}}1 {â}{{\^{a}}}1 {ș}{{\c{s}}}1 {ț}{{\c{t}}}1
           {Ă}{{\u{A}}}1 {Î}{{\^{I}}}1 {Â}{{\^{A}}}1 {Ș}{{\c{S}}}1 {Ț}{{\c{T}}}1
}

\lstdefinelanguage{TypeScript}{
  keywords={const, let, var, function, return, if, else, for, while, do, switch, case, break, continue, new, this, typeof, instanceof, in, of, class, extends, implements, interface, export, import, from, default, async, await, try, catch, finally, throw, type, enum, null, undefined, true, false, void, never, any, string, number, boolean, object, symbol, bigint},
  sensitive=true,
  morecomment=[l]{//},
  morecomment=[s]{/*}{*/},
  morestring=[b]",
  morestring=[b]',
  morestring=[b]`,
  keywordstyle=\color{blue!70!black}\bfseries,
  commentstyle=\color{green!50!black}\itshape,
  stringstyle=\color{red!60!black},
}

\lstdefinelanguage{SQL}{
  keywords={CREATE, TABLE, INSERT, SELECT, UPDATE, DELETE, FROM, WHERE, AND, OR, NOT, IN, JOIN, LEFT, RIGHT, INNER, ON, AS, ORDER, BY, GROUP, HAVING, LIMIT, OFFSET, UNION, EXISTS, BETWEEN, LIKE, IS, NULL, PRIMARY, KEY, FOREIGN, REFERENCES, DEFAULT, UNIQUE, CHECK, INDEX, USING, CASCADE, FUNCTION, RETURNS, BEGIN, END, RETURN, QUERY, LANGUAGE, plpgsql, REPLACE, EXTENSION, IF, ALTER, ADD, COLUMN, POLICY, WITH, TRUE, FALSE, COALESCE, AVG, COUNT, ANY, ARRAY, INTEGER, TEXT, UUID, DATE, BOOLEAN, NUMERIC, FLOAT, DOUBLE, PRECISION, TIMESTAMPTZ, VECTOR},
  sensitive=false,
  morecomment=[l]{--},
  morecomment=[s]{/*}{*/},
  morestring=[b]',
  keywordstyle=\color{blue!70!black}\bfseries,
  commentstyle=\color{green!50!black}\itshape,
  stringstyle=\color{red!60!black},
}

% ─── Hyperref config ─────────────────────────────────────────────────────────
\hypersetup{
  colorlinks=true,
  linkcolor=black,
  citecolor=blue!70!black,
  urlcolor=blue!70!black,
  bookmarks=true,
  pdfauthor={Bardas Cristian},
  pdftitle={PrimeRent — Platformă web pentru închirieri imobiliare cu inteligență artificială},
}

% ─── Header/Footer ──────────────────────────────────────────────────────────
\pagestyle{fancy}
\fancyhf{}
\fancyhead[L]{\small\leftmark}
\fancyhead[R]{\small\thepage}
\renewcommand{\headrulewidth}{0.4pt}

% ─── Formatare capitole/secțiuni ─────────────────────────────────────────────
\titleformat{\chapter}[display]
  {\normalfont\Large\bfseries}
  {Capitolul \thechapter}{10pt}{\LARGE}
\titlespacing*{\chapter}{0pt}{-20pt}{20pt}

% ─── Simboluri check/cross ───────────────────────────────────────────────────
\newcommand{\cmark}{\ding{51}}
\newcommand{\xmark}{\ding{55}}

% =============================================================================
\begin{document}

% ─── PAGINA DE TITLU ─────────────────────────────────────────────────────────
\begin{titlepage}
\centering
\vspace*{1cm}

{\Large\textbf{Universitatea Transilvania din Brașov}}\\[0.3cm]
{\large Facultatea de Matematică și Informatică}\\[0.3cm]
{\large Domeniul Informatică}\\[2cm]

{\Huge\textbf{LUCRARE DE LICENȚĂ}}\\[1.5cm]

{\LARGE\textbf{PrimeRent — Platformă web pentru\\închirieri imobiliare pe termen scurt\\cu integrare de inteligență artificială}}\\[3cm]

\begin{minipage}[t]{0.45\textwidth}
\raggedright
{\large\textbf{Coordonator științific:}}\\
{\large [Coordonator]}
\end{minipage}
\hfill
\begin{minipage}[t]{0.45\textwidth}
\raggedleft
{\large\textbf{Absolvent:}}\\
{\large Bardas Cristian}
\end{minipage}

\vfill
{\large Brașov, 2026}
\end{titlepage}

% ─── CUPRINS ─────────────────────────────────────────────────────────────────
\tableofcontents
\newpage

% =============================================================================
% CAPITOLUL 1 — INTRODUCERE
% =============================================================================
\chapter{INTRODUCERE}

\section{Contextul general}

În ultimii douăzeci de ani, industria închirierilor imobiliare pe termen scurt a trecut printr-o schimbare de paradigmă. Ceea ce la origini părea un concept marginal — posibilitatea de a închiria un apartament sau o cameră direct de la proprietar, fără agenții imobiliare — s-a transformat într-un sector de proporții globale. Platforme precum Airbnb, Booking.com și Vrbo au rescris regulile din industria ospitalității, oferind călătorilor alternative variate, iar proprietarilor un canal de monetizare a spațiilor locative.

Cifrele confirmă amploarea acestui fenomen. Potrivit datelor raportate de Statista, piața globală de \textit{vacation rentals} a depășit pragul de 100 de miliarde de dolari ca venituri anuale în 2023. Ritmul de creștere proiectat oscilează între 3\% și 4\% pe an în perioada 2024--2028, iar baza de utilizatori activi la nivel mondial se apropie de un miliard de persoane \cite{statista2024}.

Piața românească nu face excepție de la acest trend. Orașele mari — București, Cluj-Napoca, Brașov sau Timișoara — au văzut o explozie a ofertei de proprietăți listate pe platformele internaționale. Creșterea numărului de turiști străini, extinderea infrastructurii digitale și adoptarea pe scară largă a modelului de \textit{sharing economy} alimentează această dinamică. Zonele turistice tradiționale, de la litoral la stațiunile montane, adaugă o altă dimensiune cererii.

Pe măsură ce piața a crescut, s-au modificat și așteptările utilizatorilor. Călătorul modern nu se mai rezumă la a căuta un pat; el caută o experiență personalizată, un proces de rezervare fără fricțiuni și informații relevante prezentate clar. Din ce în ce mai frecvent, dorește recomandări adaptate preferințelor sale, nu doar filtre mecanice. Gazdele, pe de altă parte, au nevoie de instrumente care să le simplifice operațiunile zilnice: de la crearea unui anunț atractiv la gestionarea calendarului de ocupare.

Tocmai aici intervine inteligența artificială ca element diferențiator. Modelele de limbaj de dimensiuni mari (LLM-urile), exemplificate de GPT-4 al OpenAI, deschid posibilități concrete: generarea automată de texte descriptive, căutare bazată pe înțelegerea limbajului natural, sinteză a recenziilor. Tehnicile de embeddings vectoriali permit compararea semantică a textelor, depășind limitările căutării clasice prin cuvinte cheie. Aceste capabilități, integrate într-o platformă de închirieri, pot ridica semnificativ calitatea experienței pentru toți participanții.


\section{Motivația alegerii temei}

Mai mulți factori au convergit în alegerea acestei teme.

\textbf{Relevanța practică.} Sectorul de închirieri pe termen scurt este unul dinamic, cu aplicabilitate imediată. Construirea unei platforme funcționale în acest domeniu dovedește capacitatea de a traduce cerințe de business reale în soluții tehnice viabile.

\textbf{Provocarea tehnică.} Un proiect de acest calibru presupune jonglarea cu o multitudine de concepte din ingineria software: de la arhitectura client-server și bazele de date relaționale, la autentificare și autorizare, procesarea plăților, gestiunea fișierelor media, internaționalizare sau design responsive. E un teren fertil pentru demonstrarea competențelor acumulate în timpul studiilor.

\textbf{Integrarea inteligenței artificiale.} Un obiectiv central al lucrării a fost să arate cum modulele AI pot fi încorporate organic într-o aplicație web, nu ca un accesoriu decorativ, ci ca un element care produce valoare tangibilă. PrimeRent integrează trei module distincte: generare automată de conținut pentru anunțuri (inclusiv pe bază de imagini), căutare semantică în limbaj natural și sumarizare inteligentă a recenziilor. Niciuna dintre platformele comerciale de referință nu expune astfel de funcționalități direct utilizatorului final.

\textbf{Stack tehnologic modern.} Tehnologiile utilizate — React cu TypeScript, Supabase, Vercel — sunt reprezentative pentru ceea ce piața IT solicită în prezent. Alegerea lor nu este întâmplătoare, ci reflectă dorința de a lucra cu unelte relevante și de a dobândi experiență transferabilă.

Nu în ultimul rând, motivația personală: pasiunea pentru dezvoltarea web full-stack și curiozitatea de a experimenta integrarea practică a AI-ului în produse cu impact direct asupra utilizatorilor.


\section{Obiectivele lucrării}

\textbf{Obiectivul principal:} proiectarea și realizarea unei platforme web complete de închirieri imobiliare pe termen scurt, cu funcționalități integrate de inteligență artificială menite să îmbunătățească experiența atât a oaspeților, cât și a gazdelor.

\textbf{Obiective secundare:}

\begin{enumerate}[leftmargin=2cm]
  \item \textbf{Arhitectură scalabilă} — construirea unei aplicații SPA (Single Page Application) cu React și TypeScript, conectată la un backend serverless prin Supabase, cu separare clară a responsabilităților.

  \item \textbf{Autentificare securizată} — implementarea autentificării prin email și Google OAuth, cu gestiunea sesiunilor și a rolurilor (oaspete, gazdă, administrator).

  \item \textbf{Căutare avansată} — filtre multiple (locație, interval de date, număr de oaspeți, facilități), combinate cu un modul de căutare semantică prin embeddings vectoriali care acceptă interogări formulate în limbaj natural.

  \item \textbf{Module de inteligență artificială:}
  \begin{itemize}
    \item generare automată de titluri și descrieri pentru proprietăți, pe baza datelor textuale și a imaginilor (AI multimodal);
    \item căutare semantică prin convertirea interogărilor utilizatorilor în vectori și compararea lor cu vectorii proprietăților stocați în baza de date;
    \item sumarizare automată a recenziilor, cu extragerea punctelor forte și a aspectelor negative.
  \end{itemize}

  \item \textbf{Flux complet de rezervare} — de la navigare și selecție, prin verificarea disponibilității și calculul prețului, la procesarea plății prin Stripe și urmărirea stării rezervării.

  \item \textbf{Experiență de utilizare modernă} — design responsive cu abordare \textit{mobile-first}, componente accesibile, hărți interactive pentru vizualizarea locațiilor, suport bilingv (română și engleză).

  \item \textbf{Management pentru gazde} — interfață dedicată pentru crearea, editarea și administrarea anunțurilor, precum și gestionarea rezervărilor primite.

  \item \textbf{Publicare pe infrastructură cloud} — deployment pe Vercel și Supabase Cloud, cu CI/CD automatizat și accesibilitate globală.
\end{enumerate}


\section{Contribuții originale}

Lucrarea aduce câteva contribuții care o diferențiază de platformele comerciale existente:

\begin{itemize}[leftmargin=1.5cm]
  \item \textbf{AI multimodal în fluxul de creare a anunțurilor.} Sistemul procesează simultan textul introdus de gazdă și fotografiile proprietății pentru a genera descrieri atractive. Această capabilitate multimodală nu se regăsește în oferta standard a platformelor existente.

  \item \textbf{Căutare semantică prin embeddings vectoriali.} În loc să forțeze utilizatorul să completeze filtre rigide, motorul de căutare acceptă descrieri în limbaj liber (de exemplu: „un apartament confortabil lângă plajă, potrivit pentru un cuplu") și returnează rezultatele cele mai relevante din punct de vedere semantic.

  \item \textbf{Sumarizare AI a recenziilor.} O proprietate populară acumulează zeci sau chiar sute de recenzii. Funcționalitatea de sumarizare extrage automat avantajele și dezavantajele, oferind o imagine de ansamblu rapidă fără a fi nevoie de lectură individuală.
\end{itemize}


\section{Structura lucrării}

Documentul este organizat în nouă capitole, completate de bibliografie și anexe.

\textbf{Capitolul 1} — cel de față — introduce contextul, motivația, obiectivele și contribuțiile originale.

\textbf{Capitolul 2} trece în revistă piața de închirieri pe termen scurt, analizează comparativ platformele existente și identifică oportunitățile de inovare pe care PrimeRent le adresează.

\textbf{Capitolul 3} prezintă detaliat fiecare tehnologie și instrument din stack-ul ales — React, TypeScript, Vite, Tailwind CSS, Supabase, OpenAI API, Leaflet, Stripe, Vercel — justificând selecția fiecăruia.

\textbf{Capitolul 4} detaliază cerințele funcționale și non-funcționale, schema bazei de date, arhitectura de ansamblu și fluxurile principale ale sistemului.

\textbf{Capitolul 5} constituie nucleul lucrării. Aici este descrisă implementarea propriu-zisă a fiecărui modul: autentificare, căutare, rezervări, management proprietăți, funcționalități AI și interfață.

\textbf{Capitolul 6} acoperă strategia de testare, infrastructura configurată (Vitest, React Testing Library) și cele 40 de teste implementate — unitare, de componentă și de integrare.

\textbf{Capitolul 7} documentează procesul de deployment pe Vercel, configurarea serviciilor cloud și practicile de CI/CD adoptate.

\textbf{Capitolul 8} servește drept manual de utilizare, ilustrând principalele fluxuri atât din perspectiva oaspetelui, cât și a gazdei.

\textbf{Capitolul 9} sintetizează concluziile, prezintă limitările actuale și trasează direcțiile de dezvoltare viitoare.


% =============================================================================
% CAPITOLUL 2 — ANALIZA DOMENIULUI
% =============================================================================
\chapter{ANALIZA DOMENIULUI ȘI A SOLUȚIILOR EXISTENTE}

\section{Piața de închirieri imobiliare pe termen scurt}

\subsection{Evoluția și dimensiunea pieței}

Piața închirierilor pe termen scurt — desemnată adesea prin termenul englezesc \textit{short-term rentals} sau \textit{vacation rentals} — nu mai seamănă cu ceea ce era acum două decenii. Conceptul de \textit{sharing economy}, prin care resursele private sunt puse la dispoziția altora contra cost, a dat naștere unui sector care astăzi rivalizează cu industria hotelieră clasică.

Rapoartele Statista din 2024 indică un volum global de aproximativ 107 miliarde de dolari, iar proiecțiile sugerează o depășire a pragului de 130 de miliarde până în 2028 \cite{statista2024}. Rata anuală compusă de creștere se situează în jurul valorii de 3,5\%, alimentată de factori structurali: expansiunea clasei de mijloc la nivel planetar, ieftinirea transportului aerian, preferința generațiilor Millennial și Z pentru experiențe autentice în detrimentul standardizării hoteliere și, bineînțeles, digitalizarea accelerată a comportamentului de consum.

La scară europeană, activitatea este concentrată în destinațiile mediteraneene — Spania, Italia, Grecia, Croația — dar și în capitalele majore. România se înscrie în acest curent: oferta de proprietăți pe platformele internaționale a crescut vizibil, cu precădere în orașele universitare și în stațiunile montane ori litorale \cite{grandview2024}.

\subsection{Dinamica utilizatorilor}

Două categorii definesc piața.

\textbf{Oaspeții} formează un grup tot mai divers. Nu mai vorbim doar despre turiștii care caută o cameră pentru un weekend. Familii cu copii preferă apartamentele complete; nomazi digitali au nevoie de internet rapid și de un birou; grupuri de prieteni caută vile private. Fiecare segment vine cu propriul set de așteptări, ceea ce pune presiune pe platforme să ofere un grad ridicat de personalizare.

\textbf{Gazdele} s-au stratificat și ele. La un capăt al spectrului se află proprietarul individual care listează o singură locuință ocazional. La celălalt, operatori profesioniști administrează portofolii de zeci sau sute de unități. Nevoile celui de-al doilea grup sunt cu totul altele: management eficient al calendarului, comunicare rapidă, optimizarea prețurilor pe baza cererii.

\subsection{Tendințe actuale}

Câteva direcții marchează evoluția recentă a industriei:

\begin{itemize}
  \item \textbf{Profesionalizarea gazdelor.} Proporția celor care tratează închirierea ca pe un business, nu ca pe un venit ocazional, crește constant.
  \item \textbf{Reglementarea pieței.} Autoritățile locale din tot mai multe orașe impun restricții: limite ale numărului de nopți pe an, licențe, taxe specifice.
  \item \textbf{Experiența mobilă.} Peste 60\% din rezervări se efectuează de pe telefon, ceea ce face ca interfața mobilă să fie un factor critic.
  \item \textbf{Inteligența artificială.} Jucătorii mari investesc masiv în AI — de la personalizarea recomandărilor la optimizarea dinamică a prețurilor.
\end{itemize}


\section{Analiza platformelor existente}

\subsection{Airbnb}

Fondată în 2008, Airbnb rămâne liderul pieței. Peste 7 milioane de anunțuri active, în mai bine de 220 de țări și regiuni. Platforma și-a construit identitatea pe autenticitate și diversitate: de la camere modeste la castele de lux.

Printre punctele tari se numără recunoașterea globală a brandului, sistemul de recenzii bidirecționale, programul AirCover (asigurare pentru gazde) și interfața mobilă matură. Slăbiciunile? Comisioanele cumulate pot ajunge la 19\% din valoarea rezervării, suportul în disputele gazdă-oaspete este adesea criticat, iar funcționalitățile AI nu sunt vizibile utilizatorului final \cite{airbnb}.

\subsection{Booking.com}

Parte din grupul Booking Holdings, platforma acoperă atât hoteluri, cât și proprietăți private, cu un inventar impresionant de peste 28 de milioane de anunțuri. Se distinge prin sistemul de filtrare și comparare bine pus la punct, transparența prețurilor și programul de loialitate Genius.

Dezavantajele includ comisioane variabile între 10 și 25\%, o interfață care poate deveni copleșitoare și lipsa elementelor de comunitate sau de legătură personală între gazdă și oaspete. Ca și la Airbnb, funcționalitățile AI orientate spre utilizator lipsesc complet \cite{booking}.

\subsection{Vrbo}

Vrbo, deținut de Expedia Group, se concentrează exclusiv pe proprietăți întregi, nu pe camere individuale. Punctul forte: o bază fidelă de utilizatori din segmentul familiilor și grupurilor. Punctul slab: prezență geografică mai limitată, interfață mai puțin modernă, absența totală a AI-ului accesibil utilizatorilor.

\subsection{Analiza comparativă}

Tabelul \ref{tab:comparatie} sintetizează diferențele relevante între platformele analizate și PrimeRent.

\begin{table}[H]
\centering
\caption{Comparație funcționalități: PrimeRent vs.\ platforme existente}
\label{tab:comparatie}
\small
\begin{tabular}{l c c c c}
\toprule
\textbf{Funcționalitate} & \textbf{Airbnb} & \textbf{Booking} & \textbf{Vrbo} & \textbf{PrimeRent} \\
\midrule
Căutare semantică AI         & \xmark & \xmark & \xmark & \cmark \\
Generare automată descrieri   & \xmark & \xmark & \xmark & \cmark \\
Sumarizare recenzii AI        & \xmark & \xmark & \xmark & \cmark \\
Hărți interactive             & \cmark & \cmark & \cmark & \cmark \\
Calendar disponibilitate      & \cmark & \cmark & \cmark & \cmark \\
Autentificare Google OAuth    & \cmark & \xmark & \xmark & \cmark \\
Suport multi-limbă            & \cmark & \cmark & \cmark & \cmark \\
Picker locație pe hartă       & \xmark & \xmark & \xmark & \cmark \\
Codebase open-source          & \xmark & \xmark & \xmark & \cmark \\
Comisioane                    & 3--16\% & 10--25\% & 8--13\% & 0\% \\
\bottomrule
\end{tabular}
\end{table}


\section{Lacune identificate și oportunități}

\subsection{Absența funcționalităților AI vizibile utilizatorului}

Platformele mari investesc în AI la nivel intern — pentru detectarea fraudelor, optimizarea algoritmilor de ranking sau stabilirea prețurilor dinamice. Dar utilizatorul final nu vede nimic din toate acestea. Nu poate căuta în limbaj natural. Nu beneficiază de un asistent care să-i redacteze anunțul. Nu primește un rezumat al recenziilor.

Aici se deschide o oportunitate reală: oferirea de funcționalități AI accesibile, transparente și utile direct în interfața pe care utilizatorul o folosește zi de zi.

\subsection{Costul și complexitatea pentru gazdele mici}

Comisioanele practicate de platformele mari erodează marja gazdelor cu portofolii mici. La aceasta se adaugă complexitatea procesului de creare a unui anunț, care poate descuraja utilizatorii fără experiență în marketing sau copywriting.

\subsection{Lipsa flexibilității tehnice}

Toate platformele analizate sunt soluții proprietare (\textit{closed-source}). Nu permit personalizări, nu pot fi adaptate la specificul unei piețe regionale, nu se integrează ușor cu sisteme locale de administrare a proprietăților.


\section{Justificarea dezvoltării PrimeRent}

Analiza efectuată relevă un spațiu clar pentru inovare. PrimeRent se poziționează pe mai multe axe:

\begin{itemize}
  \item \textbf{AI nativ, la suprafață.} Inteligența artificială nu rămâne ascunsă în backend, ci devine un instrument activ, accesibil atât gazdelor (la redactarea anunțurilor), cât și oaspeților (la căutare și evaluarea proprietăților).
  \item \textbf{Tehnologie modernă, open-source.} Stack-ul ales permite deployment independent, fără legare de o platformă comercială centralizată.
  \item \textbf{Experiență simplificată.} Accent pe claritate vizuală, design responsive și reducerea fricțiunilor din fluxurile cheie.
  \item \textbf{Bariere reduse pentru gazde.} Funcționalitățile AI de generare a conținutului elimină nevoia de competențe de redactare pentru a crea un anunț convingător.
\end{itemize}


% =============================================================================
% CAPITOLUL 3 — TEHNOLOGII ȘI INSTRUMENTE UTILIZATE
% =============================================================================
\chapter{TEHNOLOGII ȘI INSTRUMENTE UTILIZATE}

Alegerea stack-ului tehnologic este una dintre cele mai consequente decizii dintr-un proiect software. Ea influențează viteza de dezvoltare, calitatea codului, performanța aplicației și ușurința cu care produsul poate fi extins pe termen lung. Acest capitol detaliază fiecare componentă tehnologică a platformei PrimeRent, motivând selecția făcută.

\section{Tehnologii Frontend}

\subsection{React 18}

React este o bibliotecă JavaScript open-source creată de Meta, destinată construirii interfețelor utilizator \cite{react}. De la apariția sa în 2013, a schimbat modul în care se gândesc aplicațiile web, prin introducerea conceptului de Virtual DOM și a paradigmei bazate pe componente.

Ideea centrală este \textbf{componentizarea}: interfața e decompusă în piese reutilizabile, independente, ușor de testat. Fiecare componentă își encapsulează structura (JSX), logica și starea. În PrimeRent, acest principiu se reflectă în organizarea codului: componente de pagină (\texttt{HomePage}, \texttt{ExplorePage}), de layout (\texttt{Navbar1}, \texttt{Footer}) și componente atomice (\texttt{Button}, \texttt{Card}).

\textbf{JSX} — extensia de sintaxă care permite scrierea structurii HTML direct în JavaScript — este transformat la compilare în apeluri \texttt{React.createElement()}.

\textbf{Hook-urile}, introduse în React 16.8, permit folosirea stării și a efectelor secundare în componente funcționale. Proiectul folosește extensiv hook-urile standard (\texttt{useState}, \texttt{useEffect}, \texttt{useRef}, \texttt{useCallback}) și definește un hook personalizat — \texttt{useFavorite} — care encapsulează logica de toggle a favoritelor.

Versiunea 18 aduce îmbunătățiri notabile: Concurrent Features pentru procesarea paralelă a actualizărilor de UI, Automatic Batching pentru optimizarea re-render-urilor și Suspense pentru încărcarea lazy a componentelor. PrimeRent folosește \texttt{lazy()} și \texttt{Suspense} pentru harta proprietăților individuale, amânând încărcarea Leaflet până când este efectiv necesar.

\subsection{TypeScript}

TypeScript, dezvoltat de Microsoft, adaugă tipizare statică opțională peste JavaScript \cite{typescript}. Codul TypeScript este compilat (sau, mai exact, transpilat) în JavaScript standard înainte de execuție, ceea ce înseamnă că orice browser care rulează JavaScript poate executa rezultatul.

Ce aduce concret în PrimeRent? Mai multe lucruri, fiecare cu impact direct asupra calității codului. În primul rând, \textbf{detectarea erorilor la compile-time}: greșelile de tip — accesarea unei proprietăți inexistente pe un obiect, pasarea unui argument de tip nepotrivit la o funcție — sunt semnalate de editor înainte ca aplicația să fie rulată. Aceasta reduce drastic numărul de erori descoperite abia în producție.

În al doilea rând, \textbf{autocomplete-ul și IntelliSense} devin precise. Editoarele precum VS Code pot oferi sugestii corecte de completare bazate pe tipurile definite, accelerând vizibil viteza de scriere a codului.

În al treilea rând, \textbf{interfețele servesc ca documentație vie} a structurii datelor. Dacă un dezvoltator nou deschide proiectul, nu trebuie să citească documentație separată pentru a înțelege ce câmpuri are un obiect de tip \texttt{Property} — interfața TypeScript descrie totul explicit:

\begin{lstlisting}[language=TypeScript, caption={Interfața Property din PrimeRent}]
export interface Property {
  id: string;
  title: string;
  description: string;
  city: string;
  address: string;
  lat: number;
  lng: number;
  price_per_night: number;
  max_guests: number;
  main_image: string;
  host_full_name: string;
  avg_rating: number;
  bedrooms: number;
  beds: number;
  bathrooms: number;
  is_active: boolean;
  property_images?: { url: string }[];
}
\end{lstlisting}

Un ultim avantaj important: \textbf{refactorizarea sigură}. Când se modifică structura unui tip, compilatorul TypeScript identifică automat toate locurile din cod care trebuie actualizate, eliminând riscul de a uita un loc și de a introduce un bug subtil.

Configurația proiectului (\texttt{strict: true}, \texttt{noUnusedLocals}, \texttt{noUnusedParameters}) impune cel mai înalt nivel de verificare disponibil. Opțiunea \texttt{strict} activează simultan mai multe verificări riguroase: \texttt{strictNullChecks} (valorile \texttt{null} și \texttt{undefined} trebuie tratate explicit), \texttt{strictFunctionTypes} (verificare strictă a semnăturilor funcțiilor) și altele.

\subsection{Vite}

Vite (pronunțat „veet'', din franceză — „repede'') este un build tool modern creat de Evan You, același dezvoltator din spatele Vue.js \cite{vite}. Abordarea sa diferă radical de bundlerele tradiționale precum Webpack sau Create React App.

În \textbf{development}, Vite nu bundlează întregul cod al aplicației. În schimb, folosește \textbf{modulele ES native} ale browserului și un server de development bazat pe \textbf{esbuild} (scris în Go, deci extrem de rapid). Browserul solicită modulele pe măsură ce are nevoie de ele, iar Vite le transformă individual, la cerere. Rezultatul: un server de development care pornește aproape instantaneu, indiferent de dimensiunea proiectului. Hot Module Replacement (HMR) reflectă modificările din cod în browser în milisecunde, fără a reîncărca pagina și fără a pierde starea aplicației.

În \textbf{producție}, Vite recurge la Rollup pentru bundling final, aplicând optimizări agresive: tree-shaking (eliminarea codului mort), code splitting (împărțirea în chunk-uri încărcate la cerere) și minificare.

Configurarea Vite pentru PrimeRent include aliasuri de import, plugin-ul React și plugin-ul Tailwind CSS:

\begin{lstlisting}[language=TypeScript, caption={Configurarea Vite (vite.config.ts)}]
export default defineConfig({
  plugins: [react(), tailwindcss()],
  resolve: {
    alias: { "@": path.resolve(__dirname, "./src") },
  },
  test: {
    globals: true,
    environment: "jsdom",
    setupFiles: "./src/test/setup.ts",
  },
})
\end{lstlisting}

Aliasul \texttt{@} permite importuri scurte (\texttt{@/components/Button} în loc de \texttt{../../components/Button}), simplificând navigarea între module.

\subsection{Tailwind CSS 4}

Tailwind CSS adoptă filozofia \textit{utility-first}: în loc de a scrie clase semantice precum \texttt{card} sau \texttt{primary-button} și a le stiliza în fișiere CSS separate, dezvoltatorul aplică clase utilitare de nivel jos direct în markup \cite{tailwind}. Un exemplu concret din PrimeRent:

\begin{lstlisting}[language=TypeScript, caption={Exemplu de stilizare Tailwind CSS în JSX}]
<div className="grid grid-cols-1 sm:grid-cols-2
              lg:grid-cols-3 gap-6">
  {properties.map((p) => (
    <div className="bg-white rounded-xl shadow-md
                  hover:shadow-lg transition-all
                  duration-300 overflow-hidden">
      {/* ... */}
    </div>
  ))}
</div>
\end{lstlisting}

Abordarea poate părea verbosă la prima vedere, dar elimină complet nevoia de a naviga între fișiere JSX și CSS. Stilul rămâne colocalizat cu structura componentei, iar la citirea codului se înțelege imediat cum arată fiecare element.

Design-ul responsive se realizează prin clase prefixate: \texttt{sm:} pentru ecrane $\geq 640$px, \texttt{md:} pentru $\geq 768$px, \texttt{lg:} pentru $\geq 1024$px. Astfel, \texttt{grid-cols-1 sm:grid-cols-2 lg:grid-cols-3} înseamnă: o coloană pe telefon, două pe tabletă, trei pe desktop.

\textbf{Tailwind CSS 4} aduce mai multe îmbunătățiri față de versiunea anterioară: configurare fără fișierul \texttt{tailwind.config.js} (totul se face prin plugin-ul Vite), performanță superioară la build, suport nativ pentru container queries și cascade layers CSS. Proiectul utilizează \texttt{@tailwindcss/vite} ca plugin integrat.

Un avantaj major al Tailwind este eliminarea automată a CSS-ului neutilizat la build (tree-shaking). Un proiect cu \textit{toate} clasele Tailwind ar genera un fișier CSS de peste 3MB; după purge, CSS-ul final al PrimeRent este de ordinul câtorva zeci de kilobytes.

\subsection{Radix UI și Shadcn/ui}

\textbf{Radix UI} furnizează componente React \textit{headless} — fără stilizare, dar cu comportamente complexe de accesibilitate conforme WAI-ARIA: management al focus-ului, navigare cu tastatura, atribute ARIA corecte \cite{radix}.

\textbf{Shadcn/ui} construiește peste Radix UI, adăugând stilizare cu Tailwind CSS \cite{shadcn}. Componentele sunt copiate direct în proiect, oferind control complet asupra codului. PrimeRent folosește: \texttt{Button}, \texttt{Card}, \texttt{Dialog}, \texttt{Calendar}, \texttt{Popover}, \texttt{Sheet} și altele.

\subsection{React Router v7}

React Router este biblioteca standard pentru rutarea în aplicațiile React de tip SPA \cite{reactrouter}. Versiunea 7 aduce îmbunătățiri semnificative: tipizare TypeScript completă, integrare mai strânsă cu Vite și un API simplificat.

Rutele aplicației PrimeRent acoperă toate paginile necesare celor trei tipuri de utilizatori:

\begin{table}[H]
\centering
\caption{Rutele definite în aplicația PrimeRent}
\small
\begin{tabular}{l l l}
\toprule
\textbf{Rută} & \textbf{Componenta} & \textbf{Acces} \\
\midrule
\texttt{/}                    & \texttt{HomePage}           & Public \\
\texttt{/auth}                & \texttt{AuthPage}            & Public \\
\texttt{/explore}             & \texttt{ExplorePage}         & Public \\
\texttt{/explore/:id}         & \texttt{PropertyPage}        & Public \\
\texttt{/bookings}            & \texttt{BookingsPage}        & Autentificat \\
\texttt{/favorites}           & \texttt{FavoritesPage}       & Autentificat \\
\texttt{/profile}             & \texttt{ProfilePage}         & Autentificat \\
\texttt{/host/dashboard}      & \texttt{HostDashboardPage}   & Host \\
\texttt{/host/reservations}   & \texttt{HostReservationsPage} & Host \\
\texttt{/host/create-listing} & \texttt{CreateListingPage}   & Host \\
\texttt{/host/edit-listing/:id} & \texttt{CreateListingPage}  & Host \\
\texttt{/host/manage-listings} & \texttt{ManageListingsPage} & Host \\
\texttt{/admin/embeddings}    & \texttt{AdminEmbeddingsPage} & Admin \\
\bottomrule
\end{tabular}
\end{table}

Ruta \texttt{/explore/:id} este dinamică — parametrul \texttt{:id} este extras prin hook-ul \texttt{useParams()} și folosit pentru a interoga proprietatea corespunzătoare din baza de date.

\subsection{Leaflet și react-leaflet}

Leaflet este cea mai populară bibliotecă open-source pentru hărți interactive \cite{leaflet}. Utilizează datele OpenStreetMap. În PrimeRent servește două scopuri: afișarea proprietăților ca markere cu preț pe hartă (\texttt{PropertiesMap}) și selectarea locației exacte prin click la crearea unui anunț (\texttt{CreateListingPage}).

O incompatibilitate cu Vite (căile relative ale iconițelor de marker) a fost rezolvată prin înlocuirea iconurilor implicite cu icoane personalizate.

\subsection{Internaționalizare cu i18next}

Frameworkul i18next \cite{i18next}, combinat cu \texttt{react-i18next}, asigură suportul bilingv. Detectarea limbii urmează ordinea: localStorage, setările browserului, cu fallback pe engleză. Traducerile sunt organizate pe domenii în fișiere JSON separate.


\section{Backend — Supabase}

\subsection{Prezentare generală}

Supabase se prezintă ca alternativa open-source la Firebase, cu o diferență esențială: în loc de o bază de date NoSQL, este construit pe \textbf{PostgreSQL} \cite{supabase}.

Dintr-un singur proiect, Supabase oferă: bază de date relațională cu API REST auto-generat, sistem de autentificare, stocare de fișiere, funcții serverless Edge și capabilități de timp real.

\subsection{PostgreSQL — Baza de date relațională}

PostgreSQL este motorul de baze de date central al platformei. Alegerea unui sistem relațional față de alternativele NoSQL (precum Firebase Firestore sau MongoDB) se justifică prin mai mulți factori legați de natura datelor din aplicație \cite{postgresql}.

În primul rând, datele PrimeRent au \textbf{relații clare și structurate}: proprietăți care aparțin unor gazde, rezervări care leagă oaspeți de proprietăți, recenzii asociate atât proprietăților cât și autorilor lor. Un model relațional reflectă natural aceste legături prin chei externe și constrângeri de integritate referențială.

În al doilea rând, \textbf{tranzacțiile ACID} (Atomicity, Consistency, Isolation, Durability) sunt esențiale pentru operațiuni critice precum crearea unei rezervări — care trebuie să verifice disponibilitatea, să insereze înregistrarea și să actualizeze starea, toate într-o singură operațiune atomică. Dacă oricare din pași eșuează, toate se anulează.

În al treilea rând, PostgreSQL permite \textbf{interogări complexe cu JOIN-uri multiple}, necesare pentru a returna, de exemplu, o proprietate cu numele gazdei, imaginile asociate și rating-ul mediu calculat din recenzii — totul într-un singur query.

Supabase expune baza de date PostgreSQL printr-un \textbf{API REST auto-generat} (bazat pe PostgREST), accesibil din clientul JavaScript:

\begin{lstlisting}[language=TypeScript, caption={Interogare Supabase cu relații}]
const { data } = await supabase
  .from('properties')
  .select('*, profiles(full_name), property_images(url)')
  .eq('id', propertyId)
  .single();
\end{lstlisting}

Notația \texttt{profiles(full\_name)} realizează automat un JOIN cu tabela \texttt{profiles} prin cheia externă, fără a fi nevoie de SQL explicit. Pe lângă API-ul REST, proiectul definește și \textbf{funcții RPC} (Remote Procedure Call) — funcții PostgreSQL scrise în PL/pgSQL, apelate din client. Cea mai importantă este \texttt{search\_properties}, care execută un query complex cu filtrare pe disponibilitate, locație, capacitate și facilități.

\subsection{Row Level Security (RLS)}

Row Level Security este unul dintre cele mai puternice mecanisme de securitate din PostgreSQL, și totodată unul dintre principalele motive pentru care Supabase este o alegere solidă pentru aplicații cu date sensibile \cite{supabase-rls}.

Principiul este simplu: se definesc politici declarative care specifică ce rânduri poate accesa fiecare utilizator. De exemplu, un oaspete poate vedea doar propriile rezervări, iar o gazdă poate modifica doar propriile proprietăți. Aceste reguli sunt aplicate \textit{la nivelul bazei de date}, nu în codul aplicației — ceea ce înseamnă că, chiar dacă cineva reușește să ocolească frontend-ul și să trimită cereri direct la API, politicile RLS blochează accesul neautorizat.

Mecanismul funcționează transparent: Supabase atașează automat token-ul JWT al utilizatorului autentificat la fiecare cerere, iar politicile RLS evaluează expresia \texttt{auth.uid()} pentru a determina identitatea solicitantului. Nu e nevoie de nicio linie de cod suplimentară în frontend — securitatea este inherentă schemei.

Exemplu de politică RLS din PrimeRent:

\begin{lstlisting}[language=SQL, caption={Exemplu politică RLS pentru tabela bookings}]
-- Oaspetele vede doar propriile rezervari
CREATE POLICY "Guests can view own bookings"
ON bookings FOR SELECT
USING (auth.uid() = guest_id);
\end{lstlisting}

\subsection{Supabase Auth}

Subsistemul de autentificare al Supabase oferă un set complet de capabilități, suportând mai mulți provideri de identitate. PrimeRent utilizează două:

\begin{itemize}
  \item \textbf{Email și parolă} — înregistrare clasică cu confirmare prin email, autentificare, recuperare parolă. Parolele sunt hash-uite cu bcrypt pe server; aplicația nu le stochează niciodată în clear-text.
  \item \textbf{Google OAuth} — autentificare cu un singur click prin contul Google. La prima utilizare, Supabase creează automat un utilizator în tabela \texttt{auth.users}, iar trigger-ul PostgreSQL inserează un profil corespunzător în tabela \texttt{profiles}.
\end{itemize}

Sesiunile sunt gestionate automat prin \texttt{localStorage}: la autentificare, Supabase stochează token-ul JWT și îl reîmprospătează automat înainte de expirare. În aplicație, starea de autentificare este monitorizată global în \texttt{App.tsx} prin \texttt{supabase.auth.onAuthStateChange()}, care notifică toate componentele la orice schimbare (login, logout, expirare sesiune).

\subsection{Supabase Storage}

Stocare de fișiere compatibilă cu S3, organizată în bucket-uri. Proiectul folosește un bucket \texttt{properties} pentru imaginile proprietăților. Fluxul: selecție locală, upload cu nume unic, obținere URL public.

\subsection{Edge Functions}

Funcții serverless în TypeScript/Deno, rulate la periferia rețelei \cite{supabase-edge}. Patru Edge Functions sunt definite: \texttt{generate-listing}, \texttt{smart-search}, \texttt{summarize-reviews} și \texttt{generate-embeddings}.

\subsection{pgvector — Extensia pentru căutare semantică}

Extensia pgvector adaugă suport pentru vectori n-dimensionali în PostgreSQL \cite{pgvector}. Fiecare proprietate are un vector de 1536 dimensiuni, generat de modelul \texttt{text-embedding-ada-002}. La o interogare semantică, query-ul utilizatorului este convertit în același spațiu vectorial, iar similaritatea cosinus identifică proprietățile cele mai relevante.

\subsection{Stripe}

Lider global în procesarea plăților online \cite{stripe}. Integrarea se face prin Edge Function: cheia secretă Stripe nu este niciodată expusă în browser. Utilizatorul este redirecționat la pagina securizată Stripe pentru introducerea datelor cardului.


\section{Inteligență Artificială — OpenAI API}

\subsection{GPT-4 pentru generare și sumarizare text}

\textbf{GPT-4} (Generative Pre-trained Transformer 4) este un model de limbaj de mari dimensiuni (LLM) dezvoltat de OpenAI \cite{openai}. Antrenat pe un corpus vast de text, GPT-4 demonstrează capacitatea de a înțelege și genera text în limbaj natural cu o calitate remarcabilă, urmând instrucțiuni complexe prin tehnica \textbf{prompt engineering} — construirea atentă a instrucțiunilor care ghidează comportamentul modelului.

În PrimeRent, GPT-4 este utilizat în două Edge Functions cu scopuri distincte:

\textbf{1. \texttt{generate-listing}} — Primește ca input un set de date structurate despre proprietate: titlul curent, adresa, numărul de camere, prețul, facilitățile selectate. Opțional, primește și prima imagine a proprietății codificată în base64. Generează un titlu optimizat (maxim 10 cuvinte) și o descriere atractivă de 150--200 de cuvinte. Suportul pentru imagini (multimodal) permite modelului să menționeze detalii vizuale din fotografii — o fereastră panoramică, o grădină înflorită — care nu sunt explicit introduse în formular.

\textbf{2. \texttt{summarize-reviews}} — Primește lista textelor recenziilor unei proprietăți și generează un obiect JSON structurat cu trei câmpuri: o frază de sinteză generală, o listă de puncte forte (\textit{pros}) și o listă de aspecte negative (\textit{cons}). Tehnica \texttt{response\_format: \{type: "json\_object"\}} forțează modelul să returneze exclusiv JSON valid, eliminând necesitatea de a parsa text liber și reducând riscul de erori.

Alegerea modelului \texttt{gpt-4o-mini} pentru sumarizare (în loc de \texttt{gpt-4o} complet) reflectă un compromis deliberat între cost și calitate: sumarizarea nu necesită același nivel de rafinament ca generarea de descrieri, iar modelul mini este semnificativ mai ieftin per token.

\subsection{text-embedding-ada-002 pentru căutare semantică}

\textbf{text-embedding-ada-002} este un model OpenAI specializat în generarea de \textbf{embeddings} — reprezentări vectoriale dense ale textului într-un spațiu matematic cu 1536 dimensiuni \cite{openai-embeddings}. Ideea fundamentală: texte cu înțeles similar sunt reprezentate prin vectori apropiați în acest spațiu (distanță cosinus mică), indiferent de vocabularul exact folosit.

De exemplu, textele „un apartament confortabil lângă plajă'' și „cozy beachfront flat'' vor genera vectori apropiați, chiar dacă nu au niciun cuvânt comun.

Fluxul complet al căutării semantice în PrimeRent implică doi pași:

\textbf{Pasul 1 — Indexare (offline).} Edge Function \texttt{generate-embeddings} procesează proprietățile din baza de date în batch-uri de 50. Pentru fiecare, concatenează titlul cu descrierea, generează embedding-ul prin API-ul OpenAI și îl stochează în coloana \texttt{embedding} (de tip \texttt{vector(1536)}) din tabela \texttt{properties}. Acest pas se rulează periodic din panoul de administrare.

\textbf{Pasul 2 — Căutare (real-time).} Când un utilizator introduce un query în limbaj natural, Edge Function \texttt{smart-search} generează embedding-ul query-ului cu același model, apoi execută în PostgreSQL o interogare de tip \textit{nearest-neighbor} prin extensia pgvector. Operatorul \texttt{<=>} calculează distanța cosinus între vectorul query-ului și toți vectorii proprietăților, returnând cele mai apropiate (cele mai relevante semantic).

Această arhitectură permite interogări precum „un apartament modern cu vedere la mare, ideal pentru muncă remote'' să returneze proprietăți relevante chiar dacă nicio descriere nu conține literal acele cuvinte.


\section{Deployment și DevOps}

\subsection{Vercel}

Vercel este platforma cloud pentru hosting-ul aplicațiilor frontend \cite{vercel}. Oferă CI/CD automat: fiecare push pe main declanșează build și deployment. Fișierele statice sunt distribuite pe o rețea CDN globală. HTTPS-ul este asigurat automat.

\subsection{Git și GitHub}

Git servește ca sistem de control al versiunilor; GitHub găzduiește repository-ul și facilitează integrarea cu Vercel.


\section{Instrumente de dezvoltare}

\textbf{ESLint} — analiză statică a codului JavaScript/TypeScript, cu reguli specifice React Hooks.

\textbf{Vitest și React Testing Library} — framework de testare nativ pentru Vite, combinat cu utilitare pentru testarea componentelor React prin simularea comportamentului real al utilizatorului. Detalii în Capitolul 6.

\textbf{Visual Studio Code} — editorul de cod ales pentru ecosistemul bogat de extensii și integrarea nativă cu TypeScript.


% =============================================================================
% CAPITOLUL 4 — ANALIZA CERINȚELOR ȘI PROIECTARE
% =============================================================================
\chapter{ANALIZA CERINȚELOR ȘI PROIECTARE}

\section{Cerințe funcționale}

Cerințele funcționale descriu ce trebuie să facă sistemul. PrimeRent deservește trei categorii de utilizatori — oaspeți, gazde și administratori — fiecare cu nevoi proprii.

\subsection{Cerințe pentru oaspeți}

\textbf{CF-G01 — Înregistrare și autentificare.} Crearea contului trebuie să fie un proces lipsit de fricțiuni. Platforma oferă posibilitatea creării contului clasic prin perechea email/parolă (inclusiv un flux sigur de resetare a parolei prin link pe email), sau o alternativă mult mai rapidă: autentificarea Google OAuth2 (Single Sign-On). La prima autentificare cu Google, Supabase capturează automat datele publice (nume, adresă de email, avatar) și creează o înregistrare în baza de date cu rolul implicit de \texttt{guest}.

\textbf{CF-G02 — Căutare avansată cu filtre complexe.} Motorul de căutare trebuie să permită combinarea flexibilă a filtrelor. Utilizatorii pot specifica orașul dorit, intervalul exact al șederii (check-in și check-out), numărul de oaspeți pentru care se caută cazare, și o selecție dintr-o listă de 20 de facilități posibile (ex: WiFi, Piscină, Aer condiționat, Parcare gratuită, Pet friendly). Sistemul de sortare permite aranjarea rezultatelor după prețul crescător/descrescător sau după cel mai bun rating. Toate aceste filtre sunt aplicate asincron, fără reîncărcarea completă a paginii.

\textbf{CF-G03 — Căutare semantică prin Inteligență Artificială.} O funcționalitate inovatoare care elimină necesitatea utilizării rigide a filtrelor. Oaspetele poate descrie în limbaj natural, exact așa cum i-ar vorbi unui agent de turism uman, ce fel de locuință caută (ex: „Vreau o cabană izolată la munte, perfectă pentru două persoane, cu șemineu și vedere la pădure''). Edge Function-ul \texttt{smart-search} preia textul, generează vectorul matematic, îl compară cu toate descrierile din baza de date și extrage automat numărul de oaspeți (2) și setările de mediu (munte). Rezultatele sunt sortate pur și simplu după relevanța semantică.

\textbf{CF-G04 — Vizualizare pe hartă interactivă.} Explorarea rezultatelor nu se face doar sub formă de listă (grid de carduri). Printr-un switch simplu, utilizatorul trece în modul hartă (implementat cu Leaflet). Harta afișează markere cu prețul per noapte la locațiile exacte ale proprietăților. La darea de zoom sau pan (mișcarea hărții), utilizatorul poate identifica vizual distribuția geografică a cazărilor. Un click pe marker deschide un popup cu o fotografie mică, titlul și rating-ul, care acționează ca un link către pagina detaliată a proprietății respective.

\textbf{CF-G05 — Pagina de prezentare a proprietății.} Acesta este locul unde oaspetele ia decizia de închiriere. Pagina trebuie să ofere o experiență vizuală bogată: o galerie foto superioară (imaginea principală extinsă, însoțită de un grid lateral cu alte patru poze), titlu, rating mediu bazat pe numărul total de recenzii, informații complete despre gazdă (cu poză de profil și nume), o hartă mică statică centrată pe locația proprietății, o secțiune descriptivă text, o listă cu toate facilitățile incluse și secțiunea completă de recenzii anterioare.

\textbf{CF-G06 — Sumarizare AI a recenziilor lungi.} Citirea a 50 de recenzii pentru o singură proprietate consumă timp prețios. Oaspetele are la dispoziție un buton „Sumarizează cu AI''. GPT-4o-mini procesează în fundal absolut toate textele lăsate de oaspeții anteriori și livrează în câteva secunde un scurt paragraf narativ de concluzie („Cei mai mulți oaspeți au apreciat curățenia impecabilă...''), urmat de o listă cu \textit{Puncte Forte} (ex: Comunicare excelentă, Vedere superbă) și o listă, critică, cu \textit{Aspecte Negative} (ex: Zgomot stradal dimineața, Presiune mică la duș).

\textbf{CF-G07 — Fluxul de rezervare cu verificare în timp real a disponibilității.} Panoul de rezervare domină partea dreaptă a paginii proprietății pe desktop (sau este lipit în partea de jos pe ecranele mobile). Acesta afișează un calendar în care datele deja rezervate de alți oaspeți sunt complet dezactivate, imposibil de selectat. Oaspetele selectează data de check-in și data de check-out. Sistemul calculează \textit{instantaneu} pe măsură ce selecția este făcută: numărul de nopți, înmulțite cu tariful pe noapte, plus o taxă fixă de curățenie, afișând un total transparent, fără taxe ascunse la pașii următori.

\textbf{CF-G08 — Securitatea plăților prin Stripe.} Nicio aplicație modernă nu trebuie să proceseze direct datele bancare ale clienților. La apăsarea butonului de plată, datele rezervării sunt trimise sigur pe un Edge Function care creează o sesiune Stripe, iar utilizatorul este redirecționat în afara platformei PrimeRent, pe pagina Stripe, protejată la cele mai înalte standarde PCI-DSS. Doar la finalizarea cu succes a plății, utilizatorul revine în PrimeRent iar baza de date confirmă rezervarea.

\textbf{CF-G09 — Gestionarea transparentă a rezervărilor proprii.} În pagina „My Bookings'', oaspetele poate vedea întregul istoric al călătoriilor sale. Fiecare rezervare are un card care arată detaliile proprietății, datele, prețul total achitat și, foarte important, starea curentă prin coduri de culoare (Verde pentru \textit{Confirmed}, Galben pentru \textit{Pending}, Roșu pentru \textit{Cancelled}, Albastru pentru \textit{Completed}). De aici, oaspetele poate solicita anularea gratuită a unei rezervări confirmate, cu condiția ca check-in-ul să nu aibă loc în următoarele 24 de ore.

\textbf{CF-G10 — Lăsarea recenziilor.} Feedback-ul este esențial pentru funcționarea unei economii colaborative (peer-to-peer). După data de check-out, statusul rezervării devine automat \textit{Completed}, moment în care devine vizibil butonul de recenzie. Utilizatorul alege un rating general folosind un control vizual (între 1 și 5 stele) și lasă un scurt comentariu text. Odată trimisă, recenzia nu mai poate fi editată sau ștearsă, și va contribui la calcularea noii medii vizibile pentru toți utilizatorii viitori.

\textbf{CF-G11 — Listele de favorite.} Atunci când oaspetele doar explorează și planifică o vacanță viitoare, poate apela la iconița în formă de inimă atașată fiecărui card de proprietate. Apăsarea inimii (care devine roșie) o salvează asincron în baza de date. Pagina „Favorites'' oferă acces rapid, oricând, la proprietățile selectate, chiar dacă utilizatorul schimbă dispozitivul de pe care se autentifică.

\textbf{CF-G12 — Gestionarea profilului personal.} Pagina de setări permite utilizatorului să își personalizeze prezența pe platformă: încărcarea unei noi imagini de avatar (stocată într-un bucket separat), editarea numelui complet, adăugarea unui număr de telefon valid și a datei nașterii. Aceste date sunt absolut necesare pentru a se putea transforma ulterior din oaspete simplu în gazdă.

\textbf{CF-G13 — Upgrade la gazdă.} Trecerea de la rolul \texttt{guest} la \texttt{host}, din pagina de profil.

\subsection{Cerințe pentru gazde}

\subsection{Cerințe pentru gazde}

Transformarea unui simplu cont de oaspete (guest) într-un cont de gazdă (host) deblochează o serie de instrumente de administrare.

\textbf{CF-H01 — Tabloul de bord (Dashboard) dedicat gazdei.} Acesta este punctul central de control, accesibil dintr-un buton special adăugat în bara de navigare odată ce rolul a fost schimbat. Dashboard-ul oferă o imagine de ansamblu imediată prin afișarea de statistici sumare (număr total de proprietăți listate, număr de rezervări active, venituri estimate) și legături rapide către cele trei acțiuni esențiale: adăugarea unui nou anunț, modificarea anunțurilor existente și panoul de administrare a rezervărilor primite.

\textbf{CF-H02 — Fluxul complet de creare a unui nou anunț.} Procesul de listare a unei proprietăți trebuie să fie simplu, intuitiv și împărțit logic în secțiuni. O gazdă trebuie să completeze informații de bază (titlu captivant, descriere detaliată a spațiului și regulilor casei), detalii tehnice (număr de camere, paturi, băi, capacitatea maximă de oaspeți) și un preț competitiv per noapte. 

O atenție specială este acordată locației: în loc să solicite gazdei coordonate numerice GPS care ar crea confuzie, formularul include o \textbf{hartă interactivă bazată pe Leaflet}. Gazda trebuie doar să identifice vizual locația pe hartă și să facă un simplu click pe clădirea respectivă; coordonatele de latitudine și longitudine sunt prelucrate automat și introduse în formular, invizibil pentru utilizator. Ulterior, poate introduce și un șir text reprezentând adresa formală. 

Pentru partea vizuală, o proprietate necesită imagini de calitate. Interfața oferă un panou de încărcare de tip drag-and-drop, unde utilizatorul poate selecta mai multe imagini simultan. Acestea sunt redate vizual instantaneu în browser, cu posibilitatea eliminării celor greșite înainte de salvare. 

În final, gazda trebuie să selecteze exact ce facilități oferă spațiul său (de la Aer condiționat la Pătuț de copil) dintr-o rețea predefinită de peste 20 de opțiuni distincte.

\textbf{CF-H03 — Generare asistată AI pentru titlu și descriere.} Pentru multe persoane, redactarea unui text publicitar atractiv este o sarcină descurajantă. Funcționalitatea \textit{Magic Generate} rezolvă acest impediment direct din formularul de creare. Când gazda apasă butonul asistat, platforma colectează invizibil toate datele tehnice introduse până în acel moment (tipul de proprietate, prețul, locația de pe hartă, facilitățile marcate) plus \textbf{prima fotografie} pe care a încărcat-o. Această cerere combinată (multimodală) este expediată către modelul de Inteligență Artificială GPT-4. Inteligența analizează fotografia (poate recunoaște o fereastră mare cu vedere la munte sau un șemineu luxos) și construiește un titlu de maxim 10 cuvinte extrem de atractiv, alături de o descriere fluentă și convingătoare, subliniind aspectele cele mai puternice ale casei. Gazda primește rezultatul și îl poate edita sau publica direct, scurtând cu zeci de minute procesul de listare.

\textbf{CF-H04 — Editarea flexibilă a unui anunț existent.} Realitatea pieței impune modificări constante (schimbări de prețuri, adăugări de noi dotări). Pagina de gestionare permite reîncărcarea formularului de creare cu datele deja existente în baza de date. Gazda poate adăuga noi poze, poate șterge pozele vechi prin click direct pe iconița de ștergere aferentă previzualizării sau își poate schimba descrierea cu una nouă asistată de AI. Totul se validează și se sincronizează prin Supabase Storage fără reîncărcarea vizuală a întregii ferestre.

\textbf{CF-H05 — Modulul de Management al Anunțurilor (Listings Manager).} O pagină sub formă de tabel curat ce prezintă toate locuințele gazdei respective. Gazda vede rapid prețul, ratingul și starea (activă/inactivă) și poate accesa dintr-un singur loc butoanele de vizualizare (cum apare pagina în ochii unui oaspete), editare sau ștergere definitivă (doar dacă respectiva locuință nu este implicată în rezervări viitoare nefinalizate).

\textbf{CF-H06 — Panoul de Management al Rezervărilor primite.} Acesta este sistemul back-office al gazdei. Toate rezervările făcute de oaspeți pe proprietățile sale sunt adunate aici. Când un oaspete achită prin Stripe, rezervarea intră în status \textit{Pending} pentru gazdă. Gazda are un termen rezonabil în care trebuie să intervină și să selecteze explicit butonul verde (\textit{Accept}) sau cel roșu (\textit{Decline}), adăugând astfel o ultimă barieră umană de verificare (ex: gazda constată o defecțiune subită a conductei de apă și e nevoită să anuleze cererea de cazare, evitând astfel un oaspete nefericit la locație). Odată acceptată, rezervarea trece în status \textit{Confirmed} pe ambele conturi.

\subsection{Cerințe pentru administratori}

\textbf{CF-A01 — Generare embeddings batch.} Pagină dedicată pentru procesarea proprietăților fără embeddings, în loturi de 50.


\section{Cerințe non-funcționale}

\textbf{CNF-01 — Performanță.} Code splitting, lazy loading Leaflet, caching browser.

\textbf{CNF-02 — Securitate.} RLS pe fiecare tabel, OAuth 2.0, variabile de mediu pentru chei API, validare input, HTTPS integral.

\textbf{CNF-03 — Disponibilitate.} Arhitectura serverless asigură scalare automată; SLA Supabase de 99,9\%.

\textbf{CNF-04 — Accesibilitate.} Componente Radix UI conforme WAI-ARIA, navigare completă cu tastatura.

\textbf{CNF-05 — Responsivitate.} Layout adaptat de la 320px la 1920px, cu breakpoint-uri Tailwind.

\textbf{CNF-06 — Internaționalizare.} Texte externalizate în fișiere de traducere, detecție automată a limbii.

\textbf{CNF-07 — Mentenabilitate.} Cod modular, tipizat complet cu TypeScript, reguli ESLint aplicate.


\section{Proiectarea bazei de date}

\subsection{Schema relațională}

Baza de date PostgreSQL cuprinde opt tabele principale. Fiecare folosește UUID ca cheie primară.

\textbf{Tabela \texttt{profiles}} — extinde sistemul de autentificare Supabase cu: rol (\texttt{guest}/\texttt{host}/\texttt{admin}), nume complet, telefon, data nașterii, avatar.

\textbf{Tabela \texttt{properties}} — stochează proprietățile cu toate atributele: titlu, descriere, locație (oraș, adresă, coordonate GPS), preț pe noapte, capacitate, imagine principală, rating mediu, status activ/inactiv și vectorul de embedding pentru căutarea semantică.

\textbf{Tabela \texttt{bookings}} — rezervările: proprietate, oaspete, date check-in/check-out, preț total, status.

\textbf{Tabela \texttt{reviews}} — recenziile: proprietate, autor, rating, comentariu.

\textbf{Tabela \texttt{amenities}} și \textbf{\texttt{property\_amenities}} — catalogul facilităților și asocierea lor cu proprietățile (relație N:M).

\textbf{Tabela \texttt{property\_images}} — URL-urile imaginilor suplimentare.

\textbf{Tabela \texttt{favorites}} — relațiile utilizator-proprietate pentru sistemul de favorite.

\subsection{Relațiile dintre tabele}

\begin{itemize}
  \item \texttt{profiles} $\rightarrow$ \texttt{properties}: 1:N (o gazdă deține mai multe proprietăți)
  \item \texttt{properties} $\rightarrow$ \texttt{bookings}: 1:N
  \item \texttt{profiles} $\rightarrow$ \texttt{bookings}: 1:N
  \item \texttt{properties} $\rightarrow$ \texttt{reviews}: 1:N
  \item \texttt{properties} $\leftrightarrow$ \texttt{amenities}: N:M prin \texttt{property\_amenities}
  \item \texttt{profiles} $\leftrightarrow$ \texttt{properties} prin \texttt{favorites}: N:M
\end{itemize}


\section{Arhitectura sistemului}

PrimeRent adoptă o arhitectură \textbf{JAMstack} (JavaScript, API, Markup). Frontend-ul este o aplicație statică servită de un CDN, iar toate operațiunile dinamice sunt delegate unor API-uri și funcții serverless.

Componentele arhitecturale: browserul utilizatorului (React SPA), Vercel CDN (fișiere statice), Supabase Cloud (PostgREST API, PostgreSQL cu pgvector, Storage, Edge Functions), OpenAI API (GPT-4 și embeddings) și Stripe API (procesare plăți).


% =============================================================================
% CAPITOLUL 5 — IMPLEMENTARE
% =============================================================================
\chapter{IMPLEMENTARE}

Acest capitol descrie procesul concret de implementare a platformei PrimeRent, prezentând deciziile tehnice luate pe parcurs, fragmentele de cod relevante și provocările întâmpinate la integrarea modulelor.

\section{Structura proiectului}

Organizarea codului sursă urmează o structură modulară cu separare clară a responsabilităților, respectând convențiile consacrate în ecosistemul React:

\begin{lstlisting}[caption={Structura directoarelor proiectului PrimeRent}]
src/
  components/          # Componente reutilizabile
    ui/                # Componente Shadcn/ui (Button, Card, Dialog...)
    Navbar1.tsx        # Bara de navigare
    Footer.tsx         # Subsol
    PropertiesMap.tsx   # Harta cu markere
    PropertyCard.tsx    # Cardul unei proprietati
  pages/               # Pagini (cate una per ruta)
    HomePage.tsx
    ExplorePage.tsx
    PropertyPage.tsx
    BookingsPage.tsx
    CreateListingPage.tsx
    ...
  hooks/               # Hook-uri custom
    useFavorite.ts
  i18n/                # Internationalizare
    ro.json            # Traduceri romana
    en.json            # Traduceri engleza
  types/               # Interfete TypeScript
    property.ts
    booking.ts
    review.ts
    amenity.ts
  lib/                 # Utilitare
    utils.ts           # cn(), formatare, etc.
    supabaseClient.ts  # Client Supabase singleton
  test/                # Configurare teste
    setup.ts
\end{lstlisting}

Principiul de bază: fiecare componentă își encapsulează logica, starea și structura vizuală. Paginile compun aceste componente în layout-uri complete. Hook-urile custom extrag logica reutilizabilă din componente.

Fișierul central \texttt{App.tsx} gestionează sesiunea de autentificare la nivel global prin \texttt{onAuthStateChange()} și definește arborele de rute cu React Router. Sesiunea curentă (obiectul \texttt{session}) este propagat prin props la componentele care necesită informații despre utilizatorul autentificat — o decizie deliberată în favoarea transparenței, în detrimentul unui Context API global care ar fi adăugat complexitate inutilă la această scară.


\section{Autentificare și gestiunea utilizatorilor}

Autentificarea folosește componenta \texttt{Auth} din \texttt{@supabase/auth-ui-react}, personalizată vizual. Suportă email/parolă și Google OAuth. La crearea unui cont, un trigger PostgreSQL inserează automat o înregistrare în \texttt{profiles} cu rolul \texttt{guest}.

Navbar-ul preia rolul utilizatorului și afișează condițional opțiunile: link-ul „Dashboard" este vizibil doar pentru gazde. Upgrade-ul la rolul de gazdă se face din pagina de profil, prin actualizarea câmpului \texttt{role} în baza de date.


\section{Pagina de explorare și căutare}

\subsection{Arhitectura filtrelor}

\texttt{ExplorePage} este, din punct de vedere al complexității stării, cea mai pretențioasă pagină din aplicație. Administrează simultan: locația căutată, numărul de oaspeți, intervalul de date (check-in / check-out), setul de facilități selectate, criteriul de sortare și modul de vizualizare (listă sau hartă).

O provocare specifică a fost gestionarea filtrului de facilități. Într-o implementare naivă, fiecare bifă pe un checkbox ar declanșa un nou fetch către baza de date — comportament costisitor și iritant pentru utilizator. Soluția adoptată folosește un mecanism de \textbf{staging}: selecția temporară a facilităților se menține într-o stare locală (\texttt{stagedAmenities}), iar transferul către starea principală (care declanșează fetch-ul) se face doar la apăsarea butonului ,,Apply Filters''.

\begin{lstlisting}[language=TypeScript, caption={Mecanismul de staging al filtrelor de facilități}]
const [selectedAmenities, setSelectedAmenities] = useState<string[]>([]);
const [stagedAmenities, setStagedAmenities] = useState<string[]>([]);

const handleOpenFilterDialog = () => {
  setStagedAmenities([...selectedAmenities]);
  setFilterDialogOpen(true);
};

const applyFilters = () => {
  setSelectedAmenities([...stagedAmenities]);
  setFilterDialogOpen(false);
};
\end{lstlisting}

\subsection{Interogarea prin RPC}

Funcția \texttt{fetchProperties}, declarată cu \texttt{useCallback} pentru stabilitate referențială, apelează funcția PostgreSQL \texttt{search\_properties} prin RPC. Merită subliniat ce face această funcție pe server: verifică disponibilitatea (exclude proprietățile cu rezervări suprapuse pe intervalul selectat), filtrează case-insensitive după oraș, verifică dacă proprietatea are capacitate suficientă pentru numărul de oaspeți, și se asigură că toate facilitățile solicitate sunt prezente — nu doar câteva.

\begin{lstlisting}[language=TypeScript, caption={Apelul RPC din ExplorePage}]
const { data, error } = await supabase.rpc('search_properties', {
  p_city: location || null,
  p_guests: guests,
  p_start_date: startDate || null,
  p_end_date: endDate || null,
  p_amenities: selectedAmenities.length > 0
    ? selectedAmenities : null,
});
\end{lstlisting}

\subsection{Sortarea locală}

Sortarea rezultatelor nu necesită un query suplimentar — se aplică pe array-ul deja returnat, direct în frontend. Această decizie reduce numărul de cereri către server și accelerează răspunsul perceput de utilizator.

Un caz special apare în modul ,,Top Rated'': proprietățile fără nicio recenzie (cu \texttt{avg\_rating} egal cu \texttt{null} sau 0) trebuie plasate la finalul listei, nu la început. Implementarea folosește o funcție de comparare care tratează explicit acest caz.

\subsection{Căutarea semantică AI}

Butonul ,,AI Search'' invocă Edge Function \texttt{smart-search}, care procesează query-ul utilizatorului în trei etape. Prima: convertirea textului în embedding vectorial cu \texttt{text-embedding-ada-002}. A doua: compararea vectorului rezultat cu toți vectorii proprietăților din baza de date, prin funcția PostgreSQL \texttt{match\_properties}. A treia, opțională: extragerea de filtre structurate din textul query-ului (locație, număr de oaspeți) folosind GPT-4o-mini, și aplicarea lor automată în interfață.

Rezultatul: utilizatorul scrie ,,cozy mountain cabin for two with a fireplace'', iar aplicația returnează proprietăți relevante semantic \textit{și} setează automat filtrul de oaspeți la 2.


\section{Pagina de proprietate și rezervare}

Pagina de detalii a unei proprietăți \texttt{PropertyPage} este structurată ca un orchestrator care lansează multiple interogări la montare. Aceasta este una dintre cele mai complexe pagini din aplicație, atât din perspectiva volumului de date afișate, cât și a interacțiunilor pe care le permite.

\subsection{Încărcarea datelor}

La accesarea unei proprietăți (ruta \texttt{/explore/:id}), componenta execută \textbf{patru interogări Supabase în paralel}:

\begin{enumerate}
  \item \textbf{Datele proprietății} — inclusiv profilul gazdei (\texttt{profiles(full\_name)}) și imaginile suplimentare (\texttt{property\_images(url)}).
  \item \textbf{Amenitățile} — prin join pe \texttt{property\_amenities} și \texttt{amenities}.
  \item \textbf{Recenziile} — cu profilul autorului (\texttt{profiles(full\_name, avatar\_url)}), ordonate descrescător cronologic.
  \item \textbf{Datele rezervate} — intervalele check-in/check-out din rezervările active, necesare pentru blocajele din calendar.
\end{enumerate}

Executarea lor în paralel (fiecare este un \texttt{await} independent) asigură un timp de încărcare optim. Datele sunt stocate în state-uri separate, fiecare cu propriul flag de loading.

\subsection{Galeria foto și lightbox}

Galeria adoptă un layout asimetric: imaginea principală ocupă jumătatea stângă, iar patru imagini secundare sunt dispuse într-un grid 2$\times$2 pe dreapta. Click pe oricare imagine deschide un lightbox fullscreen, implementat manual (fără biblioteci externe), cu navigare prin săgețile tastaturii, butoane în interfață și închidere cu \texttt{Escape}.

\subsection{Formularul de rezervare}

Panoul de rezervare include un \textbf{calendar interactiv} care blochează automat datele din intervalele deja rezervate. Utilizatorul selectează data de check-in (primul click) și data de check-out (al doilea click). Calculul prețului se actualizează dinamic conform formulei:

$$\text{Total} = \text{pret\_noapte} \times \text{numar\_nopti} + \text{taxa\_curatenie}$$

Taxa de curățenie este fixă: \$85 per rezervare. Această simplificare este adecvată pentru un MVP, iar o versiune viitoare ar permite gazdei să configureze taxa individual.

La apăsarea butonului ,,Reserve'', Edge Function-ul creează o sesiune Stripe Checkout cu detaliile rezervării precompletate. Utilizatorul este redirecționat la pagina securizată Stripe pentru introducerea datelor cardului. La plată reușită, callback-ul confirmă rezervarea în baza de date.


\section{Crearea și gestionarea anunțurilor}

Componenta \texttt{CreateListingPage} servește un dublu scop: atât crearea unui anunț nou, cât și editarea unuia existent. Diferența este determinată de prezența parametrului \texttt{:id} în rută. În mod editare, toate câmpurile formularului sunt preîncărcate cu datele din baza de date.

\subsection{Picker-ul de locație}

Una dintre funcționalitățile distinctive ale PrimeRent este selectarea locației pe hartă. În loc să ceară gazdei coordonate GPS (pe care nimeni nu le cunoaște din memorie), formularul afișează o hartă interactivă. La click pe orice punct, coordonatele sunt capturate automat prin hook-ul \texttt{useMapEvents} din react-leaflet, iar un marker este plasat vizual la locația selectată.

\begin{lstlisting}[language=TypeScript, caption={Componenta LocationPicker pentru selectarea locației}]
function LocationPicker({ onLocationSelect }) {
  useMapEvents({
    click(e) {
      onLocationSelect(e.latlng.lat, e.latlng.lng);
    },
  });
  return null;
}
\end{lstlisting}

\subsection{Upload-ul imaginilor}

Fluxul de upload al imaginilor combină previzualizare instantanee cu persistare asincronă. La selecția fișierelor, se generează previzualizări locale prin \texttt{URL.createObjectURL()} — fără a trimite nimic la server. Gazda vede imediat fotografiile și le poate reordona sau șterge.

La submit, imaginile sunt uploadate în paralel prin \texttt{Promise.all()}, fiecare primind un nume unic generat cu \texttt{crypto.randomUUID()}. URL-urile publice returnate de Supabase Storage sunt apoi salvate în tabela \texttt{property\_images}.

\subsection{Generarea AI a descrierii}

Funcționalitatea ,,Generate with AI'' colectează toate datele din formular (titlu, adresă, facilități, preț, capacitate) și prima imagine (codificată base64). Acestea sunt trimise la Edge Function \texttt{generate-listing}, care construiește un prompt structurat pentru GPT-4. Când există o imagine, mesajul devine multimodal — GPT-4 Vision analizează fotografia și poate menționa elemente vizuale (o priveliște montană, o terasă, un interior luminos) în descrierea generată.

Rezultatul JSON (titlu + descriere) este populat automat în formular, dar gazda poate edita liber textul înainte de publicare.


\section{Sistemul de rezervări}

Din partea oaspetelui, \texttt{BookingsPage} afișează rezervările cu badge-uri color-coded pentru status. Anularea actualizează starea local fără refetch complet (actualizare optimistă).

Din partea gazdei, \texttt{HostReservationsPage} preia rezervările pending prin join cu tabela \texttt{properties}. Acceptarea sau refuzarea se face prin aceeași funcție \texttt{updateStatus}.

Recenziile pot fi lăsate doar după finalizarea sejurului, printr-un dialog cu rating vizual și textarea.


\section{Funcționalitățile de inteligență artificială}

\subsection{Arhitectura Edge Functions}

Toate modulele AI sunt implementate ca \textbf{Edge Functions Supabase} — funcții serverless scrise în TypeScript, rulate în runtime-ul Deno la periferia rețelei. Arhitectura Edge Functions rezolvă o problemă fundamentală a aplicațiilor SPA: necesitatea de a apela API-uri care necesită chei secrete (OpenAI, Stripe), fără a expune aceste chei în codul JavaScript care rulează în browser.

Fiecare Edge Function urmează un tipar arhitectural comun:

\begin{lstlisting}[language=TypeScript, caption={Structura generală a unei Edge Function Supabase}]
import { corsHeaders } from '../_shared/cors.ts';

Deno.serve(async (req) => {
  // 1. Gestionare CORS preflight
  if (req.method === 'OPTIONS') {
    return new Response('ok', { headers: corsHeaders });
  }

  // 2. Parsare body
  const { param1, param2 } = await req.json();

  // 3. Apel API extern (OpenAI, Stripe)
  const result = await callExternalAPI(param1, param2);

  // 4. Returnare JSON
  return new Response(
    JSON.stringify(result),
    { headers: { ...corsHeaders, 'Content-Type': 'application/json' }}
  );
});
\end{lstlisting}

Header-ele CORS sunt partajate prin modulul \texttt{\_shared/cors.ts}, evitând duplicarea configurării în fiecare funcție.

\subsection{Generare descriere (\texttt{generate-listing})}

Funcția construiește un prompt structurat din datele proprietății: locație, capacitate, facilități, preț. Când gazda a uploadat o imagine, mesajul către GPT-4 devine \textbf{multimodal} — include atât text cât și imagine codificată base64, permițând modelului să observe și să descrie detalii vizuale din fotografie.

Parametrul \texttt{response\_format: \{type: "json\_object"\}} forțează GPT-4 să returneze exclusiv JSON valid cu structura predefinită (\texttt{title} și \texttt{description}), eliminând necesitatea de parsare a textului liber și reducând riscul de erori.

\subsection{Căutare semantică (\texttt{smart-search})}

Cea mai complexă Edge Function, cu trei pași secvențiali. \textbf{Primul}: conversia query-ului utilizatorului în vector de embedding prin \texttt{text-embedding-ada-002}. \textbf{Al doilea}: invocarea funcției PostgreSQL \texttt{match\_properties}, care compară vectorul cu toți vectorii din baza de date folosind operatorul \texttt{<=>} (distanță cosinus). Sunt returnate cele mai relevante proprietăți cu scor de similaritate peste pragul de 0.75.

\textbf{Al treilea} pas, opțional dar foarte util, folosește GPT-4o-mini pentru a extrage filtre structurate din textul query-ului. Dacă utilizatorul scrie ,,apartment in Cluj for 3 guests'', modelul extrage \texttt{\{city: "Cluj", guests: 3\}} și le aplică automat în interfață.

\subsection{Sumarizare recenzii (\texttt{summarize-reviews})}

Primește textele tuturor recenziilor unei proprietăți, concatenate într-un singur string. Prompt-ul instruiește GPT-4o-mini să returneze un obiect JSON cu trei câmpuri: un rezumat narativ general, o listă de puncte forte (\texttt{pros}) și o listă de aspecte negative (\texttt{cons}). Gazda și potențialii oaspeți obțin astfel o sinteză rapidă, fără a fi nevoiți să citească zeci de recenzii individual.

\subsection{Indexare batch (\texttt{generate-embeddings})}

Funcția preia din baza de date proprietățile care nu au încă un vector de embedding generat, le procesează în loturi de maximum 50, și salvează vectorii rezultați. Rularea repetată din panoul de administrare (\texttt{/admin/embeddings}) permite acoperirea progresivă a tuturor proprietăților.


\section{Internaționalizare}

Suportul bilingv (română și engleză) este implementat prin frameworkul \texttt{i18next}, configurat cu detecție automată a limbii. Ordinea de detecție: valoarea din \texttt{localStorage} (setată la ultima comutare manuală), apoi preferința browserului (\texttt{navigator.language}), cu fallback pe engleză.

Comutarea manuală se face prin componenta \texttt{LanguageSwitcher} din navbar, care apelează \texttt{i18n.changeLanguage()} și re-renderizează aplicația. Toate textele vizibile sunt externalizate în fișierele \texttt{ro.json} și \texttt{en.json}, organizate pe domenii: \texttt{navbar}, \texttt{home}, \texttt{bookings}, \texttt{explore}, etc.

Interpolarea cu parametri numerici (de exemplu, \texttt{\{\{count\}\} proprietăți disponibile}) este suportată nativ, permițând pluralizarea corectă în ambele limbi.


\section{Componente UI reutilizabile}

Sistemul de componente Shadcn/ui oferă elemente atomice configurabile prin variante predefinite, gestionate de biblioteca \texttt{class-variance-authority} (CVA). CVA permite definirea de ,,variante'' (primary, destructive, outline) și ,,dimensiuni'' (sm, default, lg) pentru fiecare componentă, fără clase CSS condiționale dispersate prin cod.

Funcția utilitară \texttt{cn()} este omniprezentă în proiect. Combină două biblioteci:

\begin{lstlisting}[language=TypeScript, caption={Funcția utilitară cn()}]
import { clsx, type ClassValue } from "clsx";
import { twMerge } from "tailwind-merge";

export function cn(...inputs: ClassValue[]) {
  return twMerge(clsx(inputs));
}
\end{lstlisting}

\texttt{clsx} convertește obiecte condiționale în string-uri de clase CSS. \texttt{twMerge} rezolvă conflictele Tailwind: dacă o componentă primește simultan \texttt{px-4} (din definiție) și \texttt{px-6} (din props), rezultatul corect este \texttt{px-6}, nu ambele. Fără \texttt{twMerge}, ambele clase ar fi aplicate, cu rezultat imprevizibil.


% =============================================================================
% CAPITOLUL 6 — TESTARE
% =============================================================================
\chapter{TESTARE ȘI VALIDARE}

\section{Strategia de testare}

PrimeRent adoptă o abordare stratificată, inspirată din piramida de testare. La bază: teste unitare — rapide, izolate, care verifică funcții pure. La nivel intermediar: teste de componentă și de integrare, care renderizează componente React complete cu dependențe înlocuite prin mock-uri. Testele End-to-End, care ar simula un utilizator real în browser, sunt planificate pentru versiuni viitoare.

\begin{table}[H]
\centering
\caption{Distribuția testelor implementate}
\label{tab:teste}
\begin{tabular}{l l r}
\toprule
\textbf{Tip} & \textbf{Module testate} & \textbf{Nr. teste} \\
\midrule
Unitar & \texttt{cn()} — utilitar CSS           & 6 \\
Unitar & \texttt{getAmenityDetails} — mapare amenități & 7 \\
Unitar & Calcul preț rezervare                & 6 \\
Unitar & Disponibilitate date calendaristice   & 7 \\
Componentă & Navbar (\texttt{Navbar1})          & 4 \\
Componentă & Pagina principală (\texttt{HomePage}) & 5 \\
Integrare & Pagina explorare (\texttt{ExplorePage}) & 5 \\
\midrule
\textbf{Total} & \textbf{6 fișiere}             & \textbf{40} \\
\bottomrule
\end{tabular}
\end{table}


\section{Infrastructura de testare}

Stack-ul: \textbf{Vitest} ca runner (compatibil nativ cu Vite), \textbf{@testing-library/react} pentru renderizare, \textbf{@testing-library/jest-dom} pentru matcher-uri DOM, \textbf{jsdom} pentru simularea browser-ului \cite{vitest, testing-library}.

Configurarea în \texttt{vite.config.ts}: \texttt{globals: true} (funcțiile de test disponibile global), \texttt{environment: "jsdom"}, fișier de setup care importă matcher-ii DOM.


\section{Teste unitare}

Testele unitare verifică funcții pure, izolate de orice dependență externă. Sunt cele mai rapide și cele mai stabile din proiect.

\textbf{Funcția \texttt{cn()}} — 6 teste care acoperă toate scenariile de utilizare: merge simplu de clase, obiecte condiționale (\texttt{\{active: true\}}), eliminarea valorilor \texttt{false}/\texttt{null}/\texttt{undefined}, și cel mai important — rezolvarea conflictelor Tailwind:

\begin{lstlisting}[language=TypeScript, caption={Teste pentru funcția cn() — rezolvare conflicte}]
test('rezolva conflictele Tailwind CSS', () => {
  expect(cn('px-4', 'px-6')).toBe('px-6');
  expect(cn('text-red-500', 'text-blue-500'))
    .toBe('text-blue-500');
});

test('returneaza string gol pentru input gol', () => {
  expect(cn()).toBe('');
});
\end{lstlisting}

\textbf{Maparea amenităților} — 7 teste: verificarea amenităților specifice (Wifi returnează iconița Wifi), valoarea fallback pentru o amenitate necunoscută (returnează icon-ul generic \texttt{ShieldCheck}), și un test parametrizat care iterează prin toate cele 20 de amenități definite.

\textbf{Logica de rezervare} — 13 teste organizate în două grupuri. Primul grup testează calculul prețului: o noapte, șapte nopți, date egale (zero nopți, se așteaptă returnarea \texttt{null}), prețuri cu zecimale. Al doilea grup testează verificarea disponibilității: situația fără rezervări existente, suprapunere totală a datelor, suprapunere parțială la început sau la sfârșit.


\section{Teste de componentă}

Testele de componentă renderizează componente React complete într-un mediu jsdom simulat, cu dependențele externe (Supabase, i18next, React Router) înlocuite prin mock-uri.

\textbf{Navbar} — 4 teste. O provocare neașteptată: navbar-ul renderizează \textbf{două logo-uri} (unul pentru desktop, altul pentru mobil, afișate condițional prin CSS), ceea ce face ca \texttt{getByAltText('Logo')} să eșueze cu eroare de elemente multiple. Soluția: \texttt{getAllByAltText('Logo')} returnează un array, iar testul verifică existența cel puțin a unui element.

\textbf{HomePage} — 5 teste, necesitând mock-uri complexe:

\begin{lstlisting}[language=TypeScript, caption={Exemplu de mock pentru Supabase în teste}]
vi.mock('@/lib/supabaseClient', () => ({
  supabase: {
    from: vi.fn(() => ({
      select: vi.fn(() => ({
        order: vi.fn(() => ({
          limit: vi.fn(() => ({
            data: mockProperties, error: null,
          })),
        })),
      })),
    })),
  },
}));
\end{lstlisting}

Testele folosesc roluri ARIA și text vizibil pentru a localiza elementele — nu clase CSS interne, conform filozofiei Testing Library care promovează testarea din perspectiva utilizatorului.


\section{Teste de integrare}

\textbf{ExplorePage} — 5 teste, cele mai complexe din proiect. Componenta depinde de Supabase (fetch proprietăți, fetch amenități), i18next (traduceri), React Router (navigare) și apeluri la Edge Functions (AI Search).

O provocare specifică a fost necesitatea de a partaja variabile mock între factory-ul \texttt{vi.mock()} și corpul testului. Vitest ridică automat (\textit{hoisting}) apelurile \texttt{vi.mock()} la începutul fișierului, dar variabilele declarate cu \texttt{const} în corpul fișierului nu sunt disponibile la acel moment. Soluția oficială: \texttt{vi.hoisted()}, care declară variabile disponibile în contextul mock-ului.

\begin{lstlisting}[language=TypeScript, caption={Utilizarea vi.hoisted() pentru variabile mock partajate}]
const { mockRpc, mockInvoke } = vi.hoisted(() => ({
  mockRpc: vi.fn(),
  mockInvoke: vi.fn(),
}));

vi.mock('@/lib/supabaseClient', () => ({
  supabase: {
    rpc: mockRpc,
    functions: { invoke: mockInvoke },
    // ...
  },
}));

test('apeleaza search_properties la incarcare', async () => {
  render(<ExplorePage session={null} />);
  expect(mockRpc).toHaveBeenCalledWith(
    'search_properties', expect.any(Object)
  );
});
\end{lstlisting}


\section{Rezultate și concluzii privind testarea}

Toate cele 40 de teste trec în sub 2 secunde, ceea ce permite rularea lor la fiecare salvare (modul \texttt{watch}) fără a întrerupe fluxul de dezvoltare.

Distribuția — 65\% teste unitare, 35\% teste de componentă/integrare — reflectă piramida clasică de testare: o bază largă de teste rapide și ieftine, completată de un număr mai mic de teste care verifică interacțiunile reale dintre componente.

Experiența a confirmat valoarea investiției: pe parcursul refactorizărilor, testele au detectat trei regresii care altfel ar fi ajuns în producție. Costul scrierii lor (aproximativ 15\% din timpul total de dezvoltare) a fost recuperat din evitarea depanărilor ulterioare.


% =============================================================================
% CAPITOLUL 7 — DEPLOYMENT ȘI CONFIGURARE
% =============================================================================
\chapter{DEPLOYMENT ȘI CONFIGURARE}

\section{Arhitectura de deployment}

Frontend-ul, un SPA compilat ca fișiere statice, este livrat prin CDN-ul Vercel. Toate operațiunile dinamice trec prin Supabase. Avantaje: performanță ridicată (servire din cache CDN), scalare automată, suprafață de atac minimă, cost zero pe planurile gratuite.

\section{Build-ul aplicației}

Comanda \texttt{npm run build} execută două etape: verificarea tipurilor TypeScript (\texttt{tsc -b}) și bundling-ul cu Vite/Rollup. Tree-shaking elimină codul neutilizat. Hash-urile în numele fișierelor asigură cache busting. CSS-ul final, după purge-ul Tailwind, este de ordinul kilobytes.

\section{Deployment pe Vercel}

Conectarea repository-ului GitHub activează CI/CD automat: fiecare push pe \texttt{main} actualizează producția; push-urile pe alte branch-uri generează preview deployments cu URL unic. Configurarea SPA routing (redirecționare către \texttt{index.html}) este detectată automat.

\section{Variabile de mediu}

Aplicația necesită două variabile prefixate \texttt{VITE\_}: URL-ul și cheia anonimă Supabase. Cheia anonimă este publică prin design — securitatea este asigurată de RLS. Cheile secrete (OpenAI, Stripe) sunt stocate exclusiv ca Supabase Secrets, accesibile doar din Edge Functions.

\section{Configurarea Supabase}

Proiectul Supabase a fost creat în regiunea \texttt{eu-central-1}. RLS activat pe toate tabelele. Google OAuth configurat cu Client ID/Secret din Google Cloud Console. Extensia pgvector activată, cu index HNSW pentru similaritate rapidă.


% =============================================================================
% CAPITOLUL 8 — MANUAL DE UTILIZARE
% =============================================================================
\chapter{MANUAL DE UTILIZARE}

\section{Accesarea platformei și prima interacțiune}

Aplicația PrimeRent a fost construită exclusiv ca platformă web, accesibilă din orice browser modern (Chrome, Safari, Firefox, Edge), fără necesitatea descărcării sau instalării vreunui software adițional din App Store sau Google Play. Aceasta urmează abordarea „Mobile-First'', ceea ce înseamnă că interfața, butoanele și meniurile se vor rearanja natural pentru a oferi o experiență perfectă atât pe ecrane masive de 1920 de pixeli, cât și pe telefoane mobile cu ecrane subțiri de 320 de pixeli lățime.

\section{Structura și navigarea din Pagina Principală (HomePage)}

La introducerea adresei web, primul ecran (denumit adesea „Hero Section'') îl întâmpină pe utilizator cu un titlu cu o scurtă animație a cuvintelor, alături de cel mai important element: \textbf{Bara rapidă de căutare}. Utilizatorul poate introduce de aici imediat orașul de destinație și este trimis automat către rezultatele pre-filtrate din pagina principală de explorare.

Imediat sub ecranul introductiv, aplicația expune statistici live ale platformei menite să creeze un sentiment de încredere (număr total de oaspeți activi, proprietăți disponibile, recenzii autentice).
Derulând mai jos, utilizatorul va întâlni două secțiuni dinamice vitale:
\begin{itemize}
    \item \textbf{Destinații Populare (Popular Destinations):} Un carusel curat, care prezintă cele mai căutate orașe din baza de date, însoțite de fotografii atrăgătoare și o etichetă ce semnalizează câte case sunt listate momentan în acel oraș (ex: „Cluj-Napoca — 12 properties'').
    \item \textbf{Proprietăți Recomandate (Top Rated Properties):} O zonă selectă unde sistemul expune primele trei cele mai bune proprietăți de pe întreaga platformă, filtrând strict după scorul mediu al recenziilor lăsate anterior.
\end{itemize}

Partea superioară este dominată permanent de \textbf{Bara de navigare (Navbar)}. Aceasta este responsivă și inteligentă. Pe un ecran mare, ea afișează toate butoanele aliniate orizontal; pe telefonul mobil, acestea se retrag curat sub un singur buton (meniul de tip „hamburger'') pentru a nu ascunde ecranul principal. 
Bara va afișa informații diferite în funcție de cine se uită la ea: dacă ești nelogat, vezi doar opțiunile „Explore'' și „Login''; dacă ești oaspete, vor apărea „Favorites'' și „My Bookings''; iar dacă ești gazdă înregistrată, va trona opțiunea „Dashboard''.

\section{Gestiunea contului (Autentificare și Profil)}

PrimeRent suportă două rute complet diferite de autentificare pentru confortul maxim.
\begin{itemize}
    \item \textbf{Autentificarea tradițională (Email și Parolă):} Utilizatorul introduce o adresă validă de email și setează o parolă. Va trebui să acceseze propriul căsuța de email pentru a confirma proprietatea asupra adresei apăsând linkul unic primit de la serverele Supabase, abia apoi având permisiunea de a intra pe platformă.
    \item \textbf{Autentificarea rapidă Google:} Utilizatorul apasă un singur buton marcat cu sigla Google. Va fi dus pe fereastra sigură Google unde confirmă contul dorit. Sistemul revine automat și asociază instant datele, fiind gata de folosire în sub 2 secunde.
\end{itemize}

Indiferent de metodă, prima intrare este tratată ca rol de \textbf{Oaspete (Guest)}.
Dacă o persoană dorește să posteze propria casă spre închiriere, trebuie să acceseze din colțul dreapta-sus, meniul sub formă de poză de profil (Avatar), și să dea click pe \textbf{„Profile''}. Aici, utilizatorul își poate completa restul datelor de identificare (inclusiv data reală a nașterii pentru siguranța vârstei legale de contractare) și apasă butonul \textbf{„Upgrade to Host''}. Procesul de conversie se realizează asincron. De menționat că trecerea în modul gazdă nu anulează drepturile de oaspete (o gazdă poate să-și continue călătoriile folosind același cont).

\section{Fluxul Oaspetelui: Explorarea Avansată și Căutarea AI}

La apăsarea butonului principal \textbf{„Explore Destinations''}, oaspetele intră în centrul decizional al PrimeRent. Această pagină oferă două mari tipuri de explorare:

\textbf{1. Căutarea clasică prin panoul de filtre.} 
În partea superioară stă un bloc solid cu câmpuri clare: Destinație, Calendar (Selectarea datei plecării și a sosirii prin două click-uri consecutive) și Număr de oaspeți. Dacă se apasă butonul \textbf{„Filters''}, un dialog lateral (Sheet) culisează și expune o grilă cu toate cele 20 de facilități posibile (WiFi, Parcare, etc.). Utilizatorul bifează ce dorește și apasă „Apply''. Din acel moment, serverul filtrează masiv și expune un grid de carduri elegante doar cu casele disponibile \textit{strict în acea perioadă de timp}. Proprietățile indisponibile nici măcar nu vor fi returnate pentru a nu aglomera vizual. Grid-ul poate fi sortat rapid folosind meniul vertical „Sort By'' din colțul dreapta.

\textbf{2. Butonul magic „AI Smart Search''.} 
Un câmp mare domină centrul cu textul \textit{Describe your ideal place in natural language}. Oaspetele poate scrie propoziții extrem de specifice.
\textit{Exemplu de utilizare:} Oaspetele scrie: „Caut neapărat o casă în Brașov pentru 4 persoane, care să aibă vedere la munte, să accepte câini și să aibă șemineu pentru iarnă''. 
Platforma nu va căuta mecanic cuvintele „munte'' sau „șemineu'', ci va înțelege sensul cererii (spațiu montan, pet-friendly, facilități încălzire de iarnă). Fără nicio altă interacțiune din partea oaspetelui, filtrele din bara de sus se vor actualiza magic pentru „Brasov'' și „4 Guests'', iar pe ecran vor fi returnate rezultatele matematic cele mai apropiate de contextul montan solicitat.

Orice card din rezultate prezintă o \textbf{inimioară} în colț. Dacă se face click pe ea, casa este salvată direct în secțiunea \textbf{„Favorites''} pentru a nu fi pierdută ulterior, creând o experiență perfectă de window-shopping pentru vacanța de vis. Pe mobil, utilizatorul poate folosi un buton flotant din colțul stânga jos pentru a comuta din vizualizarea cu grid direct într-o vizualizare de hartă aeriană uriașă, care poate fi mișcată și vizitată cu degetele.

\section{Fluxul Oaspetelui: Pagina Proprietății, Generarea Sumarului AI și Plata}

Făcând click pe o proprietate din căutare, oaspetele deschide o vizualizare individuală detaliată (ruta \texttt{/explore/:id}). Ecranul este împărțit în două pe monitoarele mari.

La stânga, conținutul: Galerie uriașă asimetrică ce poate fi deschisă în ecran complet apăsând orice poză, date tehnice detaliate și lista exhaustivă de facilități disponibile sau tăiate (indisponibile). Imediat sub ele e harta clădirii și lista lungă de recenzii anterioare.
\textbf{Cum se folosește sumarizarea AI a recenziilor:} Imediat deasupra comentariilor tronează un card albastru intitulat \textit{AI Review Summary}. La apăsarea butonului din interior, utilizatorul trebuie să aștepte în medie 5-8 secunde, timp în care aplicația conectează direct modelul OpenAI. Când se termină, cardul se extinde și livrează un raport perfect curat, punând în evidență dacă locuința are sau nu probleme cronice (ex: „Mulți oaspeți s-au plâns de zgomotul de la stradă, dar laudă priveliștea'').

La dreapta, plutește constant panoul interactiv de rezervare.
\textbf{Cum se face plata efectivă:}
1. Oaspetele dă click pe calendar. Zilele cu fundal roșu sau gri cu linie tăiată înseamnă că locuința e ocupată de altcineva.
2. Oaspetele selectează prima zi liberă, apoi selectează ultima zi. 
3. Instantaneu, panoul calculează costurile, de exemplu: \$100/noapte $\times$ 3 nopți = \$300, adaugă taxa standardizată de curățenie de \$85, și scrie cu roșu totalul: \$385.
4. Apasă \textbf{Reserve}. Într-o fracțiune de secundă se realizează tranziția sigură spre Stripe. 
5. Pe serverul Stripe, introduce detaliile cardului bancar (PrimeRent acceptă chiar și plăți de testare prin carduri puse la dispoziție de dezvoltator). Plătește și este redirecționat în siguranță spre ecranul verde \textit{Success Payment} pe site-ul PrimeRent. 

\section{Fluxul Gazdei: Adăugarea și editarea Magică cu AI a unui apartament}

Dacă utilizatorul are setat rolul de gazdă, butonul \textbf{„Dashboard''} devine inima activității sale comerciale. Acesta dă acces rapid spre secțiunea cheie \textbf{„Create New Listing''}.

Ecranul de creare este lung, curat și intuitiv, organizat pe rânduri. Prima parte necesită titlul și descrierea. Partea de mijloc necesită un click pe o hartă mare (se folosește mouse-ul pentru a așeza „pin-ul'' virtual exact pe strada corectă). După stabilirea prețului se ajunge la secțiunea Imagini.
Utilizatorul adaugă pozele de pe propriul disc; platforma nu le urcă direct online, ci construiește niște „fantome'' (previews) ale lor, astfel încât dacă una este așezată incorect, gazda poate să o șteargă printr-un click pe iconița de coș de gunoi micșorată de sub ea, fără consecințe la trafic.

\textbf{Secretul gazdei pentru atragerea oaspeților:} În loc să scrie singură texte plictisitoare, gazda trebuie doar să completeze detaliile reci, să pună prima poză și să lase secțiunile Titlu și Descriere goale. Apăsarea butonului \textbf{„Generate with AI''} va citi tot formularul plus fotografia și, în câteva secunde, va autocompleta câmpurile de titlu și descriere cu texte de marketing rafinate la nivel profesional (care includ adesea emoji-uri relevante pentru context și referințe la detalii vizuale remarcate în propria ei fotografie). 

Gazda apasă Publish, iar fotografiile trec automat prin tunelurile securizate spre sistemul Supabase Storage, fiind în câteva milisecunde vizibile public pentru toți oaspeții, fără implicarea unui administrator manual.

Orice gafă se repară din „Manage Listings'', unde gazda are la dispoziție un tabel cu absolut toate casele proprii publicate, alături de un mic creion de Edit. Apăsarea sa deschide fix formularul anterior, precompletat, pentru un proces lin de schimbare a pozelor sau actualizare de preț în funcție de anotimp.

\section{Modificarea setărilor regionale: Limba platformei}

PrimeRent vine complet internaționalizată pentru două limbi. În bara superioară, alături de imaginea profilului, stă un buton mic (ex: \textbf{RO} / \textbf{EN}). Aceasta este legătura utilizatorului cu motorul i18next. 
Apăsarea sa va re-rula invizibil fiecare etichetă vizuală din site. De exemplu, butonul „Explore Destinations'' din pagina de start va deveni fulgerător „Explorează destinații''. Această preferință este fixată pe browser-ul (telefonul) oaspetelui. Dacă o va lăsa pe limba Română, data viitoare când intră pe PrimeRent, browser-ul „va ține minte'' și o va saluta automat în Română.

Română sau Engleză, prin componenta LanguageSwitcher. Preferința se păstrează în localStorage.


% =============================================================================
% CAPITOLUL 9 — CONCLUZII
% =============================================================================
\chapter{CONCLUZII ȘI DEZVOLTĂRI VIITOARE}

\section{Rezumatul lucrării}

Lucrarea a prezentat proiectarea, implementarea și documentarea platformei PrimeRent — o aplicație web completă pentru închirierea de proprietăți pe termen scurt. Plecând de la analiza pieței și identificarea lacunelor din soluțiile existente, a fost definit un set complet de cerințe, structurat pe trei tipuri de utilizatori. Arhitectura JAMstack cu React/TypeScript și Supabase a permis dezvoltarea rapidă a unui produs funcțional, fără gestionarea infrastructurii de server.

\section{Obiective atinse}

Toate obiectivele propuse în Capitolul 1 au fost îndeplinite integral:

\begin{itemize}
  \item SPA React cu TypeScript (13 rute, strict mode)
  \item Autentificare completă (email + Google OAuth, sistem de roluri)
  \item Sistem de rezervări cu Stripe
  \item Căutare avansată (RPC PostgreSQL, 5 tipuri de filtre + sortare)
  \item Trei module AI (generare descrieri, căutare semantică, sumarizare recenzii)
  \item Hărți interactive Leaflet
  \item Internaționalizare RO/EN
  \item 40 de teste automate
  \item Deployment automatizat pe Vercel
\end{itemize}

\section{Contribuții originale}

\textbf{Căutarea semantică} depășește limitele căutării tradiționale: un utilizator poate scrie „o cabană liniștită la munte pentru o escapadă romantică" și primește rezultate relevante, chiar dacă nicio proprietate nu conține exact acele cuvinte.

\textbf{Generarea multimodală} combină text și imagini pentru a produce descrieri profesioniste, coborând bariera de intrare pentru gazdele fără experiență în copywriting.

\textbf{Sumarizarea recenziilor} comprimă zeci de opinii într-un rezumat structurat cu pro/contra, economisind timp.

\section{Dificultăți întâmpinate}

Mai multe provocări tehnice au necesitat soluții creative:

\begin{itemize}
  \item Iconițele Leaflet în Vite — căi relative incompatibile, rezolvate prin icoane personalizate.
  \item Hoisting în Vitest — variabile indisponibile în contextul \texttt{vi.mock()}, rezolvat prin \texttt{vi.hoisted()}.
  \item Securitatea cheilor API — Edge Functions ca intermediar pentru a nu expune secrete în browser.
  \item Proprietăți fără embeddings — pagina de admin pentru generare batch.
\end{itemize}

\section{Limitări actuale}

Ca produs MVP, PrimeRent are și limite: nu trimite notificări la acceptarea/anularea rezervărilor, nu oferă mesagerie directă gazdă-oaspete, anulările nu declanșează rambursări automate prin Stripe, nu acoperă teste E2E, imaginile nu sunt optimizate la upload, prețurile sunt afișate exclusiv în USD.

\section{Direcții de dezvoltare viitoare}

\textbf{Funcționalități noi:} verificare identitate (KYC), calendar de disponibilitate gestionabil manual de gazdă, prețuri dinamice bazate pe AI, recomandări personalizate, aplicație mobilă nativă (React Native).

\textbf{Îmbunătățiri tehnice:} Server-Side Rendering cu Next.js, caching avansat (TanStack Query), monitorizare erori (Sentry), audit Lighthouse.

\section{Reflecții}

Ecosistemul React modern s-a dovedit extrem de productiv. TypeScript strict previne o întreagă clasă de erori. Supabase a redus dramatic complexitatea infrastructurală, iar Row Level Security a oferit securitate declarativă, fără cod suplimentar.

Integrarea AI a demonstrat că funcționalitățile de inteligență artificială nu mai sunt exclusivitatea echipelor mari. Prin API-urile OpenAI și Edge Functions, un singur dezvoltator poate integra capabilități de nivel enterprise.

Cele 40 de teste au detectat mai multe regresii pe parcursul refactorizărilor și au oferit încrederea necesară pentru modificări substanțiale ale codului.

\section{Concluzie finală}

PrimeRent validează faptul că stack-ul React + TypeScript + Supabase + Vercel reprezintă o combinație solidă și actuală pentru aplicații web complete. Integrarea inteligenței artificiale adaugă valoare concretă și diferențiază platforma față de soluțiile tradiționale. Proiectul a atins obiectivele propuse și constituie o bază viabilă pentru un produs comercial, necesitând în principal completarea funcționalităților identificate ca limitări și validarea pe un segment real de utilizatori.


% =============================================================================
% BIBLIOGRAFIE
% =============================================================================
\begin{thebibliography}{99}

\bibitem{statista2024}
Statista Research Department, \textit{Vacation rentals — worldwide revenue forecast}, Statista, 2024. [Online]. Disponibil la: \url{https://www.statista.com/outlook/mmo/travel-tourism/vacation-rentals/worldwide}

\bibitem{grandview2024}
Grand View Research, \textit{Short-Term Rental Market Size, Share \& Trends Analysis Report 2024–2030}, 2024. [Online]. Disponibil la: \url{https://www.grandviewresearch.com/industry-analysis/short-term-rental-market}

\bibitem{airbnb}
Airbnb Engineering, \textit{Building for Scale: Lessons from Airbnb's Frontend Architecture}, Medium Engineering Blog, 2023. [Online]. Disponibil la: \url{https://medium.com/airbnb-engineering}

\bibitem{booking}
Booking Holdings, \textit{Booking.com Developer API Documentation}, 2024. [Online]. Disponibil la: \url{https://developers.booking.com/}

\bibitem{vaswani2017}
A.~Vaswani, N.~Shazeer, N.~Parmar et al., \textit{Attention Is All You Need}, in Advances in Neural Information Processing Systems (NeurIPS), vol.~30, 2017. arXiv:1706.03762.

\bibitem{brown2020}
T.~B.~Brown, B.~Mann, N.~Ryder et al., \textit{Language Models are Few-Shot Learners (GPT-3)}, arXiv:2005.14165, 2020.

\bibitem{fowler2002}
M.~Fowler, \textit{Patterns of Enterprise Application Architecture}, Addison-Wesley Professional, 2002. ISBN: 978-0321127426.

\bibitem{kleppmann2017}
M.~Kleppmann, \textit{Designing Data-Intensive Applications}, O'Reilly Media, 2017. ISBN: 978-1449373320.

\bibitem{osmani2017}
A.~Osmani, \textit{Learning JavaScript Design Patterns}, 2nd ed., O'Reilly Media, 2017. [Online]. Disponibil la: \url{https://www.patterns.dev/}

\bibitem{react}
Meta Platforms, Inc., \textit{React Documentation}, 2024. [Online]. Disponibil la: \url{https://react.dev/}

\bibitem{typescript}
Microsoft Corporation, \textit{TypeScript Documentation}, 2024. [Online]. Disponibil la: \url{https://www.typescriptlang.org/docs/}

\bibitem{vite}
E.~You \& Vite Contributors, \textit{Vite Documentation}, 2024. [Online]. Disponibil la: \url{https://vitejs.dev/guide/}

\bibitem{tailwind}
Tailwind Labs, Inc., \textit{Tailwind CSS Documentation}, 2024. [Online]. Disponibil la: \url{https://tailwindcss.com/docs/}

\bibitem{supabase}
Supabase, Inc., \textit{Supabase Documentation}, 2024. [Online]. Disponibil la: \url{https://supabase.com/docs}

\bibitem{supabase-rls}
Supabase, Inc., \textit{Row Level Security in Supabase}, 2024. [Online]. Disponibil la: \url{https://supabase.com/docs/guides/auth/row-level-security}

\bibitem{supabase-edge}
Supabase, Inc., \textit{Supabase Edge Functions with Deno}, 2024. [Online]. Disponibil la: \url{https://supabase.com/docs/guides/functions}

\bibitem{openai}
OpenAI, \textit{OpenAI API Documentation}, 2024. [Online]. Disponibil la: \url{https://platform.openai.com/docs/}

\bibitem{openai-embeddings}
OpenAI, \textit{Embeddings — text-embedding-ada-002}, 2024. [Online]. Disponibil la: \url{https://platform.openai.com/docs/guides/embeddings}

\bibitem{openai-vision}
OpenAI, \textit{GPT-4 Vision (Multimodal Inputs)}, 2024. [Online]. Disponibil la: \url{https://platform.openai.com/docs/guides/vision}

\bibitem{stripe}
Stripe, Inc., \textit{Stripe Checkout Documentation}, 2024. [Online]. Disponibil la: \url{https://stripe.com/docs/checkout}

\bibitem{vercel}
Vercel, Inc., \textit{Vercel Deployment Documentation}, 2024. [Online]. Disponibil la: \url{https://vercel.com/docs}

\bibitem{radix}
Radix UI, \textit{Radix Primitives Documentation}, 2024. [Online]. Disponibil la: \url{https://www.radix-ui.com/primitives/docs/overview/introduction}

\bibitem{shadcn}
Shadcn, \textit{shadcn/ui — Beautifully designed components}, 2024. [Online]. Disponibil la: \url{https://ui.shadcn.com/docs}

\bibitem{reactrouter}
Remix Software, Inc., \textit{React Router v7 Documentation}, 2024. [Online]. Disponibil la: \url{https://reactrouter.com/}

\bibitem{leaflet}
Leaflet Contributors, \textit{Leaflet.js — open-source JavaScript library for interactive maps}, 2024. [Online]. Disponibil la: \url{https://leafletjs.com/reference.html}

\bibitem{i18next}
i18next Contributors, \textit{i18next Documentation}, 2024. [Online]. Disponibil la: \url{https://www.i18next.com/}

\bibitem{vitest}
Vitest Contributors, \textit{Vitest — A Vite-native testing framework}, 2024. [Online]. Disponibil la: \url{https://vitest.dev/}

\bibitem{testing-library}
Testing Library Contributors, \textit{Testing Library — Simple and complete testing utilities}, 2024. [Online]. Disponibil la: \url{https://testing-library.com/docs/}

\bibitem{pgvector}
pgvector Contributors, \textit{pgvector — Open-source vector similarity search for Postgres}, 2024. [Online]. Disponibil la: \url{https://github.com/pgvector/pgvector}

\bibitem{postgresql}
PostgreSQL Global Development Group, \textit{PostgreSQL 16 Documentation}, 2024. [Online]. Disponibil la: \url{https://www.postgresql.org/docs/16/}

\bibitem{deno}
Deno Land, Inc., \textit{Deno Runtime Documentation}, 2024. [Online]. Disponibil la: \url{https://docs.deno.com/}

\bibitem{jamstack}
Netlify, \textit{The JAMstack Architecture}, 2015. [Online]. Disponibil la: \url{https://jamstack.org/}

\bibitem{owasp2021}
OWASP Foundation, \textit{OWASP Top Ten 2021}, 2021. [Online]. Disponibil la: \url{https://owasp.org/www-project-top-ten/}

\bibitem{webvitals}
Google, \textit{Web Vitals — Essential metrics for a healthy site}, 2020. [Online]. Disponibil la: \url{https://web.dev/vitals/}

\bibitem{mckinsey2023}
McKinsey \& Company, \textit{The state of AI in 2023: Generative AI's breakout year}, McKinsey Global Institute, 2023. [Online]. Disponibil la: \url{https://www.mckinsey.com/capabilities/quantumblack/our-insights/the-state-of-ai-in-2023}

\bibitem{johnson2021}
J.~Johnson, M.~Douze, H.~Jégou, \textit{Billion-Scale Similarity Search with GPUs}, IEEE Transactions on Big Data, vol.~7, no.~3, pp.~535--547, 2021.

\end{thebibliography}


% =============================================================================
% ANEXE
% =============================================================================
\appendix

\chapter{Structura completă a bazei de date}

\section{Definițiile tabelelor (SQL)}

\begin{lstlisting}[language=SQL, caption={Definirea tabelei profiles}]
CREATE TABLE profiles (
  id         UUID REFERENCES auth.users(id) PRIMARY KEY,
  full_name  TEXT,
  phone_number TEXT,
  birth_date DATE,
  avatar_url TEXT,
  role       TEXT DEFAULT 'guest'
             CHECK (role IN ('guest', 'host', 'admin')),
  created_at TIMESTAMPTZ DEFAULT now()
);
\end{lstlisting}

\begin{lstlisting}[language=SQL, caption={Definirea tabelei properties}]
CREATE TABLE properties (
  id              UUID DEFAULT gen_random_uuid() PRIMARY KEY,
  host_id         UUID REFERENCES profiles(id) ON DELETE CASCADE,
  title           TEXT NOT NULL,
  description     TEXT,
  city            TEXT NOT NULL,
  address         TEXT,
  lat             DOUBLE PRECISION,
  lng             DOUBLE PRECISION,
  price_per_night NUMERIC(10,2) NOT NULL,
  max_guests      INTEGER DEFAULT 1,
  bedrooms        INTEGER DEFAULT 1,
  beds            INTEGER DEFAULT 1,
  bathrooms       INTEGER DEFAULT 1,
  main_image      TEXT,
  is_active       BOOLEAN DEFAULT true,
  embedding       VECTOR(1536),
  created_at      TIMESTAMPTZ DEFAULT now()
);
\end{lstlisting}

\begin{lstlisting}[language=SQL, caption={Definirea tabelelor bookings și reviews}]
CREATE TABLE bookings (
  id           UUID DEFAULT gen_random_uuid() PRIMARY KEY,
  property_id  UUID REFERENCES properties(id) ON DELETE CASCADE,
  guest_id     UUID REFERENCES profiles(id),
  host_id      UUID REFERENCES profiles(id),
  check_in     DATE NOT NULL,
  check_out    DATE NOT NULL,
  guest_count  INTEGER DEFAULT 1,
  total_price  NUMERIC(10,2),
  status       TEXT DEFAULT 'pending'
               CHECK (status IN ('pending','confirmed',
                                 'cancelled','completed')),
  created_at   TIMESTAMPTZ DEFAULT now()
);

CREATE TABLE reviews (
  id          UUID DEFAULT gen_random_uuid() PRIMARY KEY,
  property_id UUID REFERENCES properties(id) ON DELETE CASCADE,
  guest_id    UUID REFERENCES profiles(id),
  rating      INTEGER NOT NULL CHECK (rating BETWEEN 1 AND 5),
  comment     TEXT,
  created_at  TIMESTAMPTZ DEFAULT now()
);
\end{lstlisting}

\begin{lstlisting}[language=SQL, caption={Definirea tabelelor de suport}]
CREATE TABLE amenities (
  id       UUID DEFAULT gen_random_uuid() PRIMARY KEY,
  name     TEXT NOT NULL UNIQUE,
  category TEXT
);

CREATE TABLE property_amenities (
  id          UUID DEFAULT gen_random_uuid() PRIMARY KEY,
  property_id UUID REFERENCES properties(id) ON DELETE CASCADE,
  amenity_id  UUID REFERENCES amenities(id) ON DELETE CASCADE,
  UNIQUE(property_id, amenity_id)
);

CREATE TABLE property_images (
  id          UUID DEFAULT gen_random_uuid() PRIMARY KEY,
  property_id UUID REFERENCES properties(id) ON DELETE CASCADE,
  url         TEXT NOT NULL,
  created_at  TIMESTAMPTZ DEFAULT now()
);

CREATE TABLE favorites (
  id          UUID DEFAULT gen_random_uuid() PRIMARY KEY,
  property_id UUID REFERENCES properties(id) ON DELETE CASCADE,
  profile_id  UUID REFERENCES profiles(id) ON DELETE CASCADE,
  created_at  TIMESTAMPTZ DEFAULT now(),
  UNIQUE(property_id, profile_id)
);
\end{lstlisting}


\chapter{Funcțiile PostgreSQL (RPC)}

\begin{lstlisting}[language=SQL, caption={Funcția search\_properties}]
CREATE OR REPLACE FUNCTION search_properties(
  p_city       TEXT    DEFAULT NULL,
  p_guests     INTEGER DEFAULT 1,
  p_start_date DATE    DEFAULT NULL,
  p_end_date   DATE    DEFAULT NULL,
  p_amenities  UUID[]  DEFAULT NULL
)
RETURNS TABLE (
  id UUID, title TEXT, description TEXT, city TEXT,
  address TEXT, lat DOUBLE PRECISION, lng DOUBLE PRECISION,
  price_per_night NUMERIC, max_guests INTEGER,
  bedrooms INTEGER, beds INTEGER, bathrooms INTEGER,
  main_image TEXT, host_full_name TEXT,
  avg_rating NUMERIC, is_active BOOLEAN
) AS $$
BEGIN
  RETURN QUERY
  SELECT p.id, p.title, p.description, p.city, p.address,
    p.lat, p.lng, p.price_per_night, p.max_guests,
    p.bedrooms, p.beds, p.bathrooms, p.main_image,
    pr.full_name AS host_full_name,
    COALESCE(AVG(r.rating), 0)::NUMERIC AS avg_rating,
    p.is_active
  FROM properties p
  LEFT JOIN profiles pr ON pr.id = p.host_id
  LEFT JOIN reviews r ON r.property_id = p.id
  WHERE p.is_active = true
    AND (p_city IS NULL
         OR LOWER(p.city) LIKE LOWER('%' || p_city || '%'))
    AND p.max_guests >= p_guests
    AND (p_start_date IS NULL OR p_end_date IS NULL
         OR p.id NOT IN (
           SELECT b.property_id FROM bookings b
           WHERE b.status NOT IN ('cancelled')
             AND b.check_in < p_end_date
             AND b.check_out > p_start_date))
    AND (p_amenities IS NULL OR (
      SELECT COUNT(*) FROM property_amenities pa
      WHERE pa.property_id = p.id
        AND pa.amenity_id = ANY(p_amenities)
    ) = array_length(p_amenities, 1))
  GROUP BY p.id, pr.full_name
  ORDER BY avg_rating DESC;
END;
$$ LANGUAGE plpgsql;
\end{lstlisting}

\begin{lstlisting}[language=SQL, caption={Funcția match\_properties (pgvector)}]
CREATE OR REPLACE FUNCTION match_properties(
  query_embedding VECTOR(1536),
  match_threshold FLOAT DEFAULT 0.75,
  match_count     INTEGER DEFAULT 20
)
RETURNS TABLE (
  id UUID, title TEXT, description TEXT,
  city TEXT, similarity FLOAT
) AS $$
BEGIN
  RETURN QUERY
  SELECT p.id, p.title, p.description, p.city,
    1 - (p.embedding <=> query_embedding) AS similarity
  FROM properties p
  WHERE p.embedding IS NOT NULL
    AND p.is_active = true
    AND 1 - (p.embedding <=> query_embedding) > match_threshold
  ORDER BY p.embedding <=> query_embedding
  LIMIT match_count;
END;
$$ LANGUAGE plpgsql;
\end{lstlisting}


\chapter{Politicile Row Level Security (RLS)}

\begin{lstlisting}[language=SQL, caption={Politici RLS pentru tabela properties}]
CREATE POLICY "Anyone can view active properties"
ON properties FOR SELECT
USING (is_active = true);

CREATE POLICY "Hosts can create properties"
ON properties FOR INSERT
WITH CHECK (
  auth.uid() = host_id
  AND EXISTS (
    SELECT 1 FROM profiles
    WHERE id = auth.uid() AND role = 'host'
  )
);

CREATE POLICY "Hosts can update their own properties"
ON properties FOR UPDATE
USING (auth.uid() = host_id);

CREATE POLICY "Hosts can delete their own properties"
ON properties FOR DELETE
USING (auth.uid() = host_id);
\end{lstlisting}

\begin{lstlisting}[language=SQL, caption={Politici RLS pentru tabela bookings}]
CREATE POLICY "Guests can view own bookings"
ON bookings FOR SELECT
USING (auth.uid() = guest_id);

CREATE POLICY "Hosts can view bookings for their properties"
ON bookings FOR SELECT
USING (auth.uid() = host_id);

CREATE POLICY "Authenticated users can create bookings"
ON bookings FOR INSERT
WITH CHECK (auth.uid() = guest_id);

CREATE POLICY "Guests can cancel own bookings"
ON bookings FOR UPDATE
USING (auth.uid() = guest_id);

CREATE POLICY "Hosts can update bookings for their properties"
ON bookings FOR UPDATE
USING (auth.uid() = host_id);
\end{lstlisting}

\begin{lstlisting}[language=SQL, caption={Politici RLS pentru tabela favorites}]
CREATE POLICY "Users can manage own favorites"
ON favorites FOR ALL
USING (auth.uid() = profile_id)
WITH CHECK (auth.uid() = profile_id);
\end{lstlisting}


\chapter{Glosarul termenilor tehnici}

\begin{longtable}{p{4cm} p{10cm}}
\toprule
\textbf{Termen} & \textbf{Definiție} \\
\midrule
\endfirsthead
\toprule
\textbf{Termen} & \textbf{Definiție} \\
\midrule
\endhead
\textbf{SPA} & Single Page Application — aplicație web care încarcă o singură pagină HTML și actualizează dinamic conținutul \\
\textbf{JAMstack} & Arhitectură bazată pe JavaScript, API-uri și Markup pregenerat \\
\textbf{BaaS} & Backend-as-a-Service — serviciu cloud cu funcționalități de backend gata de utilizat \\
\textbf{RLS} & Row Level Security — control al accesului la date la nivel de rând \\
\textbf{JWT} & JSON Web Token — standard pentru transmiterea securizată a informațiilor de autentificare \\
\textbf{Edge Function} & Funcție serverless rulată la periferia rețelei \\
\textbf{pgvector} & Extensie PostgreSQL pentru vectori și operații de similaritate \\
\textbf{Embedding} & Reprezentare vectorială a unui text, capturând semantica \\
\textbf{Cosine Similarity} & Măsură a similarității între doi vectori, bazată pe unghiul dintre ei \\
\textbf{HMR} & Hot Module Replacement — actualizarea modulelor fără refresh complet \\
\textbf{CDN} & Content Delivery Network — rețea distribuită de servere \\
\textbf{OAuth} & Protocol de autorizare prin conturi terțe \\
\textbf{CI/CD} & Continuous Integration / Continuous Deployment \\
\textbf{Tree-shaking} & Eliminarea codului neutilizat din bundle \\
\textbf{RPC} & Remote Procedure Call — apelarea funcțiilor server din client \\
\textbf{Hoisting} & Mutarea automată a declarațiilor la topul scopului \\
\textbf{Mock} & Înlocuitor al unui modul real în teste \\
\bottomrule
\end{longtable}


\end{document}
